\hypertarget{section}{%
\section{}\label{section}}

\learningobjectives{

1. \important{}
2. \important{}  
3. \important{}
4. \important{}
5. \important{}
6. \important{}
}

\hypertarget{section-1}{%
\subsection{}\label{section-1}}

\important{}

\important{}

\important{}

\important{}

\important{}

\hypertarget{section-2}{%
\subsection{}\label{section-2}}

!!! info ``\,''

\begin{verbatim}
### [ ](section08-01.md)
AI

\important{}20200
\important{}
- -   
- - 
- ### [ ](section08-02.md)

\important{}530
\important{}
- - LSTM
- - 
- ### [ ](section08-03.md)

\important{}508
\important{}
- - 
- - 
- ### [ ](section08-04.md)

\important{}3008000
\important{}
- - 
- - 
- ### [ ](section08-05.md)

\important{}
- API
- DevOps
- - 
- - 
- \important{}

### [ ](section08-02.md)

\important{}520100

\important{}
- - 
- WebGL
- - 

\important{}CesiumLOD

### [ ](section08-03.md)

\important{}2001000

\important{}
- - 
- - 
- \important{}Apache KafkaInfluxDB

### [ ](section08-04.md)

\important{}3200

\important{}
- - 
- - 
- \important{}

### [ ](section08-05.md)

\important{}

\important{}
- - 
- - 
- \important{}DevOpsAI
\end{verbatim}

\hypertarget{section-3}{%
\subsection{}\label{section-3}}

!!! note ``\,''

\begin{verbatim}
=== ""
    - \important{} + Spring Boot + Vue.js + MySQL + Redis
    - \important{}Cesium + Three.js + WebGL + Node.js + MongoDB  
    - \important{}Kafka + InfluxDB + Spark + Python + ECharts
    - \important{}TensorFlow + FastAPI + PostgreSQL + React + Docker
    - \important{}Kubernetes + Jenkins + Prometheus + Grafana + MinIO
    
=== ""
    - \important{}
    - \important{}
    - \important{}
    - \important{}
    
=== ""
    - \important{}
    - \important{}
    - \important{}
\end{verbatim}

\hypertarget{section-4}{%
\subsection{}\label{section-4}}

\begin{longtable}[]{@{}llllll@{}}
\toprule\noalign{}
\endhead
\bottomrule\noalign{}
\endlastfoot
\important{} & & & & & \\
\important{} & Web & 3D & & & \\
\important{} & 3-6 & 6-9 & 9-12 & 12-18 & 18-24 \\
\important{} & 5-8 & 8-12 & 12-15 & 15-20 & 20-30 \\
\important{} & & & & & \\
\important{} & & & & & \\
\end{longtable}

\hypertarget{section-5}{%
\subsection{}\label{section-5}}

!!! tip ``\,''

\begin{verbatim}
=== ""
    \important{}
    - Java/Python
    - WebHTML/CSS/JavaScript
    - SQL
    - Linux
    
    \important{}
    - - 
    - Web
    
=== ""
    \important{}
    - - Hadoop/Spark
    - - 
    
    \important{}
    - - 
    - === ""
    \important{}
    - - K8s/Docker
    - DevOps
    - \important{}
    - - 
    - ## 
\end{verbatim}

!!! tip ``\,''

\begin{verbatim}
=== ""
    1. \important{}
    2. \important{}
    3. \important{}
    
=== ""
    4. \important{}
    5. \important{}
    6. \important{}
    
=== ""
    7. \important{}
    8. \important{}
    9. \important{}
\end{verbatim}

\hypertarget{section-6}{%
\subsection{}\label{section-6}}

\begin{longtable}[]{@{}llll@{}}
\toprule\noalign{}
\endhead
\bottomrule\noalign{}
\endlastfoot
& & & \\
& & & \\
& & & \\
& & & \\
& & & \\
\end{longtable}

\hypertarget{section-7}{%
\subsection{}\label{section-7}}

!!! warning ``\,''

\begin{verbatim}
- \important{}
- \important{}
- \important{}
- \important{}
- \important{}
- \important{}
- \important{}
- \important{}
\end{verbatim}

\hypertarget{section-8}{%
\subsubsection{}\label{section-8}}

\begin{itemize}
\tightlist
\item
  \important{}Spring BootSpring CloudMyBatis PlusRedisRabbitMQ
\item
  \important{}Vue.jsElement UIEChartsCesiumWebGL
\item
  \important{}MySQLInfluxDBMongoDB
\item
  \important{}Apache SparkFlinkKafka
\item
  \important{}DockerKubernetesJenkins
\item
  \important{}PrometheusGrafanaELK Stack
\end{itemize}

\hypertarget{section-9}{%
\subsubsection{}\label{section-9}}

\begin{itemize}
\item
  \important{}
\item
  \important{}
\item
  \important{}CQRS
\item
  \important{}
\end{itemize}

\hypertarget{section-10}{%
\subsubsection{}\label{section-10}}

\begin{itemize}
\item
  \important{}
\item
  \important{}
\item
  \important{}
\item
  \important{}
\item
  \important{}
\end{itemize}

\hypertarget{section-11}{%
\subsection{}\label{section-11}}

\hypertarget{section-12}{%
\subsubsection{}\label{section-12}}

\begin{itemize}
\item
  \begin{itemize}
  \tightlist
  \item
  \end{itemize}
\item ~
  \hypertarget{section-13}{%
  \subsubsection{}\label{section-13}}
\item
  \begin{itemize}
  \tightlist
  \item
  \end{itemize}
\item ~
  \hypertarget{section-14}{%
  \subsubsection{}\label{section-14}}
\item
  \begin{itemize}
  \tightlist
  \item
  \end{itemize}
\item ~
  \hypertarget{section-15}{%
  \subsubsection{}\label{section-15}}
\item
  \begin{itemize}
  \tightlist
  \item
  \end{itemize}
\item ~
  \hypertarget{section-16}{%
  \subsection{}\label{section-16}}
\end{itemize}

\hypertarget{section-17}{%
\subsubsection{\texorpdfstring{\href{section08-01.md}{}}{}}\label{section-17}}

\hypertarget{section-18}{%
\subsubsection{\texorpdfstring{\href{section08-02.md}{}}{}}\label{section-18}}

\hypertarget{section-19}{%
\subsubsection{\texorpdfstring{\href{section08-03.md}{}}{}}\label{section-19}}

\hypertarget{section-20}{%
\subsubsection{\texorpdfstring{\href{section08-04.md}{}}{}}\label{section-20}}

\hypertarget{section-21}{%
\subsubsection{\texorpdfstring{\href{section08-05.md}{}}{}}\label{section-21}}

\hypertarget{section-22}{%
\subsection{}\label{section-22}}

\hypertarget{section-23}{%
\subsubsection{}\label{section-23}}

\begin{enumerate}
\def\labelenumi{\arabic{enumi}.}
\item
  \important{}
\item
  \important{}
\item
  \important{}
\item
  \important{}
\end{enumerate}

\hypertarget{section-24}{%
\subsubsection{}\label{section-24}}

\begin{enumerate}
\def\labelenumi{\arabic{enumi}.}
\item
  \important{}
\item
  \important{}
\item
  \important{}
\item
  \important{}
\end{enumerate}

\hypertarget{section-25}{%
\subsubsection{}\label{section-25}}

\begin{enumerate}
\def\labelenumi{\arabic{enumi}.}
\item
  \important{}
\item
  \important{}
\item
  \important{}
\item
  \important{}
\end{enumerate}

\hypertarget{section-26}{%
\subsection{}\label{section-26}}

{[}1{]} . {[}M{]}. : , 2022.

{[}2{]} . {[}S{]}. 2022226, 2022.

{[}3{]} , . (2021-2025){[}S{]}. 2021.

{[}4{]} . {[}J{]}. , 2021, (21): 1-3.

{[}5{]} Martin Fowler. Microservices Architecture{[}OL{]}. 2014.
https://martinfowler.com/articles/microservices.html

{[}6{]} Newman S. Building Microservices: Designing Fine-Grained
Systems{[}M{]}. 2nd ed.~Sebastopol: O'Reilly Media, 2021.

{[}7{]} Richardson C. Microservices Patterns: With Examples in
Java{[}M{]}. Shelter Island: Manning Publications, 2018.

{[}8{]} Burns B, Beda J. Kubernetes: Up and Running{[}M{]}. 2nd
ed.~Sebastopol: O'Reilly Media, 2019.

{[}9{]} Gamma E, Helm R, Johnson R, et al.~Design Patterns: Elements of
Reusable Object-Oriented Software{[}M{]}. Boston: Addison-Wesley, 1994.

{[}10{]} . {[}R{]}. : , 2021.

{[}11{]} Apache Software Foundation. Apache Kafka Documentation{[}OL{]}.
2023. https://kafka.apache.org/documentation/

{[}12{]} The Spring Team. Spring Boot Reference Documentation{[}OL{]}.
VMware, 2023. https://spring.io/projects/spring-boot

{[}13{]} Vue.js Team. Vue.js Guide{[}OL{]}. 2023.
https://vuejs.org/guide/

{[}14{]} , , . {[}J{]}. , 2020, 51(2): 132-141.

{[}15{]} , , . {[}J{]}. , 2020, 31(1): 1-10.

{[}16{]} ISO/IEC 25010:2011. Systems and software engineering ---
Systems and software Quality Requirements and Evaluation
(SQuaRE){[}S{]}. International Organization for Standardization, 2011.

{[}17{]} GB/T 36073-2018. {[}S{]}. : , 2018.

{[}18{]} IEEE Std 1471-2000. IEEE Recommended Practice for Architectural
Description of Software-Intensive Systems{[}S{]}. IEEE Computer Society,
2000.

{[}19{]} Humble J, Farley D. Continuous Delivery: Reliable Software
Releases through Build, Test, and Deployment Automation{[}M{]}. Boston:
Addison-Wesley, 2010.

{[}20{]} Kim G, Humble J, Debois P, et al.~The DevOps Handbook: How to
Create World-Class Agility, Reliability, and Security in Technology
Organizations{[}M{]}. Portland: IT Revolution Press, 2016.

\hypertarget{section-27}{%
\subsection{}\label{section-27}}

DevOps

\hypertarget{section-28}{%
\section{}\label{section-28}}

\hypertarget{section-29}{%
\subsection{8.1.1}\label{section-29}}

\hypertarget{section-30}{%
\subsubsection{8.1.1.1}\label{section-30}}

100507×24

\important{}

\begin{itemize}
\item
  \begin{itemize}
  \tightlist
  \item
    128
  \end{itemize}
\item
  15.8
\item
  300MW
\item
  1200+
\item ~
  \hypertarget{section-31}{%
  \subsubsection{8.1.1.2}\label{section-31}}
\end{itemize}

\hypertarget{section-32}{%
\paragraph{}\label{section-32}}

\important{1. }

\begin{itemize}
\item
  \begin{itemize}
  \tightlist
  \item
  \end{itemize}
\item
  \begin{itemize}
  \tightlist
  \item
  \end{itemize}
\end{itemize}

\important{2. }

\begin{itemize}
\item
  \begin{itemize}
  \tightlist
  \item
  \end{itemize}
\item
  \begin{itemize}
  \tightlist
  \item
  \end{itemize}
\end{itemize}

\important{3. }

\begin{itemize}
\item
  \begin{itemize}
  \tightlist
  \item
  \end{itemize}
\item
  \begin{itemize}
  \tightlist
  \item
  \end{itemize}
\end{itemize}

\important{4. }

\begin{itemize}
\item
  \begin{itemize}
  \tightlist
  \item
  \end{itemize}
\item
  \begin{itemize}
  \tightlist
  \item
  \end{itemize}
\end{itemize}

\important{5. }

\begin{itemize}
\item
  \begin{itemize}
  \tightlist
  \item
  \end{itemize}
\item
  \begin{itemize}
  \tightlist
  \item
  \end{itemize}
\end{itemize}

\hypertarget{section-33}{%
\paragraph{}\label{section-33}}

\important{1. }

\begin{itemize}
\tightlist
\item
  100+
\item
  ≤3
\item
  99.9\%
\item
  PB
\end{itemize}

\important{2. }

\begin{itemize}
\item
  \begin{itemize}
  \tightlist
  \item
  \end{itemize}
\item
  \begin{itemize}
  \tightlist
  \item
  \end{itemize}
\end{itemize}

\important{3. }

\begin{itemize}
\item
  \begin{itemize}
  \tightlist
  \item
  \end{itemize}
\item
  \begin{itemize}
  \tightlist
  \item
  \end{itemize}
\end{itemize}

\hypertarget{section-34}{%
\subsubsection{8.1.1.3}\label{section-34}}

\hypertarget{section-35}{%
\paragraph{}\label{section-35}}

\begin{itemize}
\item
  \begin{itemize}
  \tightlist
  \item
  \end{itemize}
\item
  \begin{itemize}
  \tightlist
  \item
  \end{itemize}
\end{itemize}

\hypertarget{section-36}{%
\paragraph{}\label{section-36}}

\begin{itemize}
\tightlist
\item
  \important{}Vue.js + Element UI + ECharts + Cesium
\item
  \important{}Spring Boot + Spring Cloud + MyBatis Plus
\item
  \important{}MySQL + Redis + InfluxDB
\item
  \important{}RabbitMQ
\item
  \important{}Spark + Flink
\item
  \important{}Docker + Kubernetes
\end{itemize}

\hypertarget{section-37}{%
\subsection{8.1.2}\label{section-37}}

\hypertarget{section-38}{%
\subsubsection{8.1.2.1}\label{section-38}}

``--''

\begin{lstlisting}

                         AI        

                             ...

\end{lstlisting}

\hypertarget{section-39}{%
\subsubsection{8.1.2.2}\label{section-39}}

\begin{lstlisting}

                        API Gateway                         
                   (Spring Cloud Gateway)                   

   User           Device         Data       
   Service        Service        Service    

   Monitor        Alert          Analysis   
   Service        Service        Service    

                    Eureka     
                    
\end{lstlisting}

\hypertarget{section-40}{%
\subsubsection{8.1.2.3}\label{section-40}}

\hypertarget{section-41}{%
\paragraph{}\label{section-41}}

\important{1.  (monitoring_points)}
\begin{lstlisting}[style=sql]
CREATE TABLE monitoring_points (
    id BIGINT PRIMARY KEY AUTO_INCREMENT,
    point_code VARCHAR(50) UNIQUE NOT NULL COMMENT '',
    point_name VARCHAR(100) NOT NULL COMMENT '',
    point_type VARCHAR(20) NOT NULL COMMENT '',
    position_x DECIMAL(10,6) COMMENT 'X',
    position_y DECIMAL(10,6) COMMENT 'Y',
    position_z DECIMAL(8,3) COMMENT 'Z',
    sensor_type VARCHAR(50) COMMENT '',
    install_date DATE COMMENT '',
    status TINYINT DEFAULT 1 COMMENT '1- 0-',
    created_time TIMESTAMP DEFAULT CURRENT_TIMESTAMP,
    updated_time TIMESTAMP DEFAULT CURRENT_TIMESTAMP ON UPDATE CURRENT_TIMESTAMP
);
\end{lstlisting}

\important{2.  (monitoring_data)}
\begin{lstlisting}[style=sql]
CREATE TABLE monitoring_data (
    id BIGINT PRIMARY KEY AUTO_INCREMENT,
    point_id BIGINT NOT NULL COMMENT 'ID',
    measure_time TIMESTAMP NOT NULL COMMENT '',
    parameter_code VARCHAR(20) NOT NULL COMMENT '',
    parameter_value DECIMAL(12,4) COMMENT '',
    parameter_unit VARCHAR(10) COMMENT '',
    data_quality TINYINT DEFAULT 1 COMMENT '1- 2- 3-',
    created_time TIMESTAMP DEFAULT CURRENT_TIMESTAMP,
    INDEX idx_point_time (point_id, measure_time),
    INDEX idx_measure_time (measure_time)
);
\end{lstlisting}

\important{3.  (alert_rules)}
\begin{lstlisting}[style=sql]
CREATE TABLE alert_rules (
    id BIGINT PRIMARY KEY AUTO_INCREMENT,
    rule_name VARCHAR(100) NOT NULL COMMENT '',
    point_type VARCHAR(20) NOT NULL COMMENT '',
    parameter_code VARCHAR(20) NOT NULL COMMENT '',
    alert_level TINYINT NOT NULL COMMENT '1- 2- 3- 4-',
    threshold_value DECIMAL(12,4) COMMENT '',
    threshold_type TINYINT COMMENT '1- 2- 3-',
    status TINYINT DEFAULT 1 COMMENT '1- 0-',
    created_time TIMESTAMP DEFAULT CURRENT_TIMESTAMP
);
\end{lstlisting}

\hypertarget{section-42}{%
\subsection{8.1.3}\label{section-42}}

\hypertarget{section-43}{%
\subsubsection{8.1.3.1}\label{section-43}}

\hypertarget{section-44}{%
\paragraph{}\label{section-44}}

\begin{lstlisting}[style=java]
@RestController
@RequestMapping("/api/data-collection")
@Slf4j
public class DataCollectionController {
    
    @Autowired
    private DataCollectionService dataCollectionService;
    
    /**
     \highlight{ 
     }/
    @PostMapping("/batch")
    public Result<Void> batchReceiveData(@RequestBody List<MonitoringDataDTO> dataList) {
        try {
            dataCollectionService.batchProcessData(dataList);
            return Result.success();
        } catch (Exception e) {
            log.error("", e);
            return Result.error("" + e.getMessage());
        }
    }
    
    /**
     \highlight{ 
     }/
    @PostMapping("/realtime")
    public Result<Void> receiveRealtimeData(@RequestBody MonitoringDataDTO data) {
        try {
            dataCollectionService.processRealtimeData(data);
            return Result.success();
        } catch (Exception e) {
            log.error("", e);
            return Result.error("" + e.getMessage());
        }
    }
}

@Service
@Transactional
public class DataCollectionServiceImpl implements DataCollectionService {
    
    @Autowired
    private MonitoringDataMapper monitoringDataMapper;
    
    @Autowired
    private DataValidationService dataValidationService;
    
    @Autowired
    private RabbitTemplate rabbitTemplate;
    
    @Override
    public void batchProcessData(List<MonitoringDataDTO> dataList) {
        // 
        List<MonitoringData> validDataList = new ArrayList<>();
        for (MonitoringDataDTO dto : dataList) {
            if (dataValidationService.validate(dto)) {
                MonitoringData data = convertToEntity(dto);
                validDataList.add(data);
            }
        }
        
        // 
        if (!validDataList.isEmpty()) {
            monitoringDataMapper.batchInsert(validDataList);
            
            // 
            for (MonitoringData data : validDataList) {
                rabbitTemplate.convertAndSend("alert.exchange", 
                    "data.received", data);
            }
        }
    }
    
    @Override
    public void processRealtimeData(MonitoringDataDTO dto) {
        // 
        if (dataValidationService.validate(dto)) {
            MonitoringData data = convertToEntity(dto);
            
            // 
            monitoringDataMapper.insert(data);
            
            // Redis
            redisTemplate.opsForValue().set(
                "realtime:data:" + data.getPointId(), 
                data, 300, TimeUnit.SECONDS);
            
            // WebSocket
            websocketService.broadcastData(data);
            
            // 
            rabbitTemplate.convertAndSend("alert.exchange", 
                "data.realtime", data);
        }
    }
}
\end{lstlisting}

\hypertarget{section-45}{%
\subsubsection{8.1.3.2}\label{section-45}}

\hypertarget{section-46}{%
\paragraph{}\label{section-46}}

\begin{lstlisting}[style=java]
@Component
@RabbitListener(queues = "alert.queue")
@Slf4j
public class AlertAnalysisEngine {
    
    @Autowired
    private AlertRuleService alertRuleService;
    
    @Autowired
    private AlertRecordService alertRecordService;
    
    @Autowired
    private NotificationService notificationService;
    
    /**
     \highlight{ 
     }/
    @RabbitHandler
    public void analyzeData(MonitoringData data) {
        try {
            // 
            List<AlertRule> rules = alertRuleService.getRulesByPointType(
                data.getPointType(), data.getParameterCode());
            
            for (AlertRule rule : rules) {
                if (checkAlertCondition(data, rule)) {
                    // 
                    AlertRecord alert = createAlertRecord(data, rule);
                    alertRecordService.save(alert);
                    
                    // 
                    notificationService.sendAlert(alert);
                    
                    log.warn(" - {}{}{}", 
                        data.getPointCode(), data.getParameterCode(), 
                        rule.getAlertLevel());
                }
            }
        } catch (Exception e) {
            log.error("", e);
        }
    }
    
    /**
     \highlight{ 
     }/
    private boolean checkAlertCondition(MonitoringData data, AlertRule rule) {
        BigDecimal value = data.getParameterValue();
        BigDecimal threshold = rule.getThresholdValue();
        
        switch (rule.getThresholdType()) {
            case 1: // 
                return value.compareTo(threshold) > 0;
            case 2: // 
                return value.compareTo(threshold) < 0;
            case 3: // 
                return checkChangeRate(data, threshold);
            default:
                return false;
        }
    }
    
    /**
     \highlight{ 
     }/
    private boolean checkChangeRate(MonitoringData data, BigDecimal threshold) {
        // 
        MonitoringData lastData = alertRuleService.getLastData(
            data.getPointId(), data.getParameterCode());
        
        if (lastData != null) {
            BigDecimal changeRate = data.getParameterValue()
                .subtract(lastData.getParameterValue())
                .divide(lastData.getParameterValue(), 4, RoundingMode.HALF_UP)
                .multiply(new BigDecimal("100"));
            
            return changeRate.abs().compareTo(threshold) > 0;
        }
        
        return false;
    }
}
\end{lstlisting}

\hypertarget{section-47}{%
\subsubsection{8.1.3.3}\label{section-47}}

\hypertarget{section-48}{%
\paragraph{}\label{section-48}}

\begin{lstlisting}
<template>
  <div class="monitoring-3d-view">
    <div id="cesiumContainer" class="cesium-container"></div>
    
    <!--  -->
    <div class="control-panel">
      <el-card>
        <h3></h3>
        <el-tree
          :data="monitoringPoints"
          :props="treeProps"
          show-checkbox
          @check="onPointCheck"
        />
      </el-card>
      
      <el-card>
        <h3></h3>
        <div v-for="alert in alerts" :key="alert.id" 
             :class="['alert-item', 'level-' + alert.level]">
          <div class="alert-title">{{ alert.pointName }}</div>
          <div class="alert-content">{{ alert.message }}</div>
          <div class="alert-time">{{ alert.createTime }}</div>
        </div>
      </el-card>
    </div>
    
    <!--  -->
    <div class="data-panel" v-if="selectedPoint">
      <el-card>
        <h3>{{ selectedPoint.pointName }}</h3>
        <div class="data-charts">
          <div ref="trendChart" class="trend-chart"></div>
        </div>
        
        <el-table :data="selectedPoint.realtimeData" size="small">
          <el-table-column prop="parameterName" label="" width="120"/>
          <el-table-column prop="value" label="" width="100"/>
          <el-table-column prop="unit" label="" width="80"/>
          <el-table-column prop="updateTime" label=""/>
        </el-table>
      </el-card>
    </div>
  </div>
</template>

<script>
import \highlight{ as Cesium from 'cesium'
import } as echarts from 'echarts'

export default {
  name: 'Monitoring3DView',
  data() {
    return {
      viewer: null,
      monitoringPoints: [],
      alerts: [],
      selectedPoint: null,
      trendChart: null,
      pointEntities: new Map(),
      treeProps: {
        children: 'children',
        label: 'name'
      }
    }
  },
  
  mounted() {
    this.initCesium()
    this.loadMonitoringPoints()
    this.initWebSocket()
    this.loadAlerts()
  },
  
  methods: {
    /**
     \highlight{ Cesium
     }/
    initCesium() {
      // Cesium
      Cesium.Ion.defaultAccessToken = 'your_cesium_access_token'
      
      // Cesium Viewer
      this.viewer = new Cesium.Viewer('cesiumContainer', {
        terrainProvider: Cesium.createWorldTerrain(),
        skyBox: new Cesium.SkyBox({
          sources: {
            positiveX: '/assets/skybox/TychoSkymapII.t3_08192x04096_80_px.jpg',
            negativeX: '/assets/skybox/TychoSkymapII.t3_08192x04096_80_mx.jpg',
            positiveY: '/assets/skybox/TychoSkymapII.t3_08192x04096_80_py.jpg',
            negativeY: '/assets/skybox/TychoSkymapII.t3_08192x04096_80_my.jpg',
            positiveZ: '/assets/skybox/TychoSkymapII.t3_08192x04096_80_pz.jpg',
            negativeZ: '/assets/skybox/TychoSkymapII.t3_08192x04096_80_mz.jpg'
          }
        })
      })
      
      // 
      this.loadDamModel()
      
      // 
      this.viewer.camera.setView({
        destination: Cesium.Cartesian3.fromDegrees(116.3974, 39.9093, 1000),
        orientation: {
          heading: Cesium.Math.toRadians(0),
          pitch: Cesium.Math.toRadians(-30),
          roll: 0.0
        }
      })
    },
    
    /**
     \highlight{ 
     }/
    async loadDamModel() {
      try {
        const damModel = await Cesium.Model.fromGltf({
          url: '/assets/models/dam_model.gltf',
          modelMatrix: Cesium.Transforms.eastNorthUpToFixedFrame(
            Cesium.Cartesian3.fromDegrees(116.3974, 39.9093, 0)
          )
        })
        
        this.viewer.scene.primitives.add(damModel)
      } catch (error) {
        console.error(':', error)
      }
    },
    
    /**
     \highlight{ 
     }/
    async loadMonitoringPoints() {
      try {
        const response = await this.$api.get('/monitoring/points')
        this.monitoringPoints = response.data
        
        // 
        this.addPointsToScene()
      } catch (error) {
        this.$message.error('')
      }
    },
    
    /**
     \highlight{ 
     }/
    addPointsToScene() {
      this.monitoringPoints.forEach(point => {
        const entity = this.viewer.entities.add({
          id: point.id,
          position: Cesium.Cartesian3.fromDegrees(
            point.longitude, point.latitude, point.elevation
          ),
          point: {
            pixelSize: 10,
            color: this.getPointColor(point.status),
            outlineColor: Cesium.Color.WHITE,
            outlineWidth: 2,
            heightReference: Cesium.HeightReference.RELATIVE_TO_GROUND
          },
          label: {
            text: point.pointName,
            font: '14pt monospace',
            style: Cesium.LabelStyle.FILL_AND_OUTLINE,
            outlineWidth: 2,
            verticalOrigin: Cesium.VerticalOrigin.BOTTOM,
            pixelOffset: new Cesium.Cartesian2(0, -30),
            heightReference: Cesium.HeightReference.RELATIVE_TO_GROUND
          }
        })
        
        this.pointEntities.set(point.id, entity)
      })
      
      // 
      this.viewer.cesiumWidget.screenSpaceEventHandler.setInputAction(
        this.onPointClick.bind(this),
        Cesium.ScreenSpaceEventType.LEFT_CLICK
      )
    },
    
    /**
     \highlight{ 
     }/
    getPointColor(status) {
      switch (status) {
        case 'normal': return Cesium.Color.GREEN
        case 'warning': return Cesium.Color.YELLOW
        case 'alert': return Cesium.Color.RED
        case 'offline': return Cesium.Color.GRAY
        default: return Cesium.Color.WHITE
      }
    },
    
    /**
     \highlight{ 
     }/
    onPointClick(event) {
      const picked = this.viewer.scene.pick(event.position)
      if (picked && picked.id) {
        const pointId = picked.id.id
        this.selectPoint(pointId)
      }
    },
    
    /**
     \highlight{ 
     }/
    async selectPoint(pointId) {
      try {
        const response = await this.$api.get(\code{/monitoring/points/${pointId}/detail})
        this.selectedPoint = response.data
        
        // 
        this.drawTrendChart()
        
        // 
        this.highlightPoint(pointId)
      } catch (error) {
        this.$message.error('')
      }
    },
    
    /**
     \highlight{ 
     }/
    drawTrendChart() {
      if (!this.selectedPoint.trendData) return
      
      this.$nextTick(() => {
        if (this.trendChart) {
          this.trendChart.dispose()
        }
        
        this.trendChart = echarts.init(this.$refs.trendChart)
        
        const option = {
          title: {
            text: '',
            left: 'center'
          },
          tooltip: {
            trigger: 'axis'
          },
          legend: {
            data: this.selectedPoint.trendData.parameters
          },
          xAxis: {
            type: 'time',
            data: this.selectedPoint.trendData.times
          },
          yAxis: {
            type: 'value'
          },
          series: this.selectedPoint.trendData.parameters.map(param => ({
            name: param.name,
            type: 'line',
            data: param.values
          }))
        }
        
        this.trendChart.setOption(option)
      })
    },
    
    /**
     \highlight{ WebSocket
     }/
    initWebSocket() {
      const ws = new WebSocket('ws://localhost:8080/ws/monitoring')
      
      ws.onmessage = (event) => {
        const data = JSON.parse(event.data)
        
        if (data.type === 'realtime_data') {
          this.updateRealtimeData(data.payload)
        } else if (data.type === 'alert') {
          this.handleAlert(data.payload)
        }
      }
      
      ws.onerror = (error) => {
        console.error('WebSocket:', error)
        this.$message.error('')
      }
    },
    
    /**
     \highlight{ 
     }/
    updateRealtimeData(data) {
      // 
      const entity = this.pointEntities.get(data.pointId)
      if (entity) {
        entity.point.color = this.getPointColor(data.status)
      }
      
      // 
      if (this.selectedPoint && this.selectedPoint.id === data.pointId) {
        this.selectedPoint.realtimeData = data.parameters
      }
    },
    
    /**
     \highlight{ 
     }/
    handleAlert(alert) {
      // 
      this.alerts.unshift(alert)
      
      // 
      const entity = this.pointEntities.get(alert.pointId)
      if (entity) {
        entity.point.color = this.getAlertColor(alert.level)
        
        // 
        this.addBlinkEffect(entity, alert.level)
      }
      
      // 
      this.$notify({
        title: \code{${alert.levelName}},
        message: \code{${alert.pointName}: ${alert.message}},
        type: 'warning',
        duration: 0
      })
    },
    
    /**
     \highlight{ 
     }/
    getAlertColor(level) {
      switch (level) {
        case 1: return Cesium.Color.BLUE
        case 2: return Cesium.Color.YELLOW
        case 3: return Cesium.Color.ORANGE
        case 4: return Cesium.Color.RED
        default: return Cesium.Color.WHITE
      }
    },
    
    /**
     \highlight{ 
     }/
    addBlinkEffect(entity, level) {
      const originalColor = entity.point.color.getValue()
      const alertColor = this.getAlertColor(level)
      
      let isVisible = true
      const blinkInterval = setInterval(() => {
        entity.point.color = isVisible ? alertColor : originalColor
        isVisible = !isVisible
      }, 500)
      
      // 5
      setTimeout(() => {
        clearInterval(blinkInterval)
        entity.point.color = alertColor
      }, 5000)
    }
  }
}
</script>

<style scoped>
.monitoring-3d-view {
  position: relative;
  width: 100%;
  height: 100vh;
}

.cesium-container {
  width: 100%;
  height: 100%;
}

.control-panel {
  position: absolute;
  top: 20px;
  left: 20px;
  width: 300px;
  max-height: 60vh;
  overflow-y: auto;
  z-index: 1000;
}

.control-panel .el-card {
  margin-bottom: 10px;
}

.data-panel {
  position: absolute;
  top: 20px;
  right: 20px;
  width: 400px;
  max-height: 80vh;
  overflow-y: auto;
  z-index: 1000;
}

.trend-chart {
  width: 100%;
  height: 300px;
}

.alert-item {
  padding: 10px;
  margin-bottom: 10px;
  border-radius: 4px;
  background-color: #f5f5f5;
}

.alert-item.level-1 { border-left: 4px solid #409eff; }
.alert-item.level-2 { border-left: 4px solid #e6a23c; }
.alert-item.level-3 { border-left: 4px solid #f56c6c; }
.alert-item.level-4 { border-left: 4px solid #ff0000; }

.alert-title {
  font-weight: bold;
  margin-bottom: 5px;
}

.alert-content {
  font-size: 14px;
  color: #666;
}

.alert-time {
  font-size: 12px;
  color: #999;
  text-align: right;
  margin-top: 5px;
}
</style>
\end{lstlisting}

\hypertarget{section-49}{%
\subsection{8.1.4}\label{section-49}}

\hypertarget{section-50}{%
\subsubsection{8.1.4.1}\label{section-50}}

\hypertarget{docker}{%
\paragraph{Docker}\label{docker}}

\begin{lstlisting}
# Dockerfile
FROM openjdk:11-jre-slim

LABEL maintainer="water-platform@example.com"

WORKDIR /app

COPY target/water-monitoring-\highlight{.jar app.jar

EXPOSE 8080

ENTRYPOINT ["java", "-jar", "app.jar"]
\end{lstlisting}

\hypertarget{docker-compose}{%
\paragraph{Docker Compose}\label{docker-compose}}

\begin{lstlisting}
version: '3.8'

services:
  # MySQL
  mysql:
    image: mysql:8.0
    container_name: water-mysql
    environment:
      MYSQL_ROOT_PASSWORD: root123
      MYSQL_DATABASE: water_monitoring
      MYSQL_USER: water_user
      MYSQL_PASSWORD: water_pass
    volumes:
      - mysql_data:/var/lib/mysql
      - ./sql:/docker-entrypoint-initdb.d
    ports:
      - "3306:3306"
    networks:
      - water-network

  # Redis
  redis:
    image: redis:6.2-alpine
    container_name: water-redis
    ports:
      - "6379:6379"
    volumes:
      - redis_data:/data
    networks:
      - water-network

  # RabbitMQ
  rabbitmq:
    image: rabbitmq:3.9-management
    container_name: water-rabbitmq
    environment:
      RABBITMQ_DEFAULT_USER: admin
      RABBITMQ_DEFAULT_PASS: admin123
    ports:
      - "5672:5672"
      - "15672:15672"
    volumes:
      - rabbitmq_data:/var/lib/rabbitmq
    networks:
      - water-network

  # 
  water-backend:
    build: .
    container_name: water-backend
    environment:
      SPRING_PROFILES_ACTIVE: docker
      MYSQL_HOST: mysql
      REDIS_HOST: redis
      RABBITMQ_HOST: rabbitmq
    ports:
      - "8080:8080"
    depends_on:
      - mysql
      - redis
      - rabbitmq
    networks:
      - water-network

  # 
  water-frontend:
    image: nginx:alpine
    container_name: water-frontend
    volumes:
      - ./dist:/usr/share/nginx/html
      - ./nginx.conf:/etc/nginx/nginx.conf
    ports:
      - "80:80"
    depends_on:
      - water-backend
    networks:
      - water-network

volumes:
  mysql_data:
  redis_data:
  rabbitmq_data:

networks:
  water-network:
    driver: bridge
\end{lstlisting}

\hypertarget{section-51}{%
\subsubsection{8.1.4.2}\label{section-51}}

\hypertarget{section-52}{%
\paragraph{}\label{section-52}}

\begin{lstlisting}
# prometheus.yml
global:
  scrape_interval: 15s

scrape_configs:
  - job_name: 'water-monitoring'
    static_configs:
      - targets: ['water-backend:8080']
    metrics_path: '/actuator/prometheus'

  - job_name: 'mysql'
    static_configs:
      - targets: ['mysql:3306']

  - job_name: 'redis'
    static_configs:
      - targets: ['redis:6379']
\end{lstlisting}

\hypertarget{section-53}{%
\paragraph{}\label{section-53}}

\begin{lstlisting}
# alert_rules.yml
groups:
  - name: water_monitoring_alerts
    rules:
      - alert: ServiceDown
        expr: up == 0
        for: 1m
        labels:
          severity: critical
        annotations:
          summary: " {{ $labels.instance }} "
          description: "{{ $labels.job }} 1"

      - alert: HighMemoryUsage
        expr: (1 - (node_memory_MemAvailable_bytes / node_memory_MemTotal_bytes)) } 100 > 80
        for: 2m
        labels:
          severity: warning
        annotations:
          summary: ""
          description: " {{ $labels.instance }}  {{ $value }}%"

      - alert: HighCPUUsage
        expr: 100 - (avg(rate(node_cpu_seconds_total{mode="idle"}[5m])) \highlight{ 100) > 80
        for: 2m
        labels:
          severity: warning
        annotations:
          summary: "CPU"
          description: " {{ $labels.instance }} CPU {{ $value }}%"

      - alert: DatabaseConnectionPoolHigh
        expr: hikaricp_connections_active / hikaricp_connections_max } 100 > 80
        for: 1m
        labels:
          severity: warning
        annotations:
          summary: ""
          description: " {{ $value }}%"
\end{lstlisting}

\hypertarget{section-54}{%
\subsection{8.1.5}\label{section-54}}

\hypertarget{section-55}{%
\subsubsection{8.1.5.1}\label{section-55}}

\begin{enumerate}
\def\labelenumi{\arabic{enumi}.}
\item
  \important{}
\item
  \important{}
\item
  \important{}Cesium
\item
  \important{}
\item
  \important{}PC
\end{enumerate}

\hypertarget{section-56}{%
\subsubsection{8.1.5.2}\label{section-56}}

\begin{longtable}[]{@{}ll@{}}
\toprule\noalign{}
\i & mportant\{} \\
\midrule\noalign{}
\endhead
\bottomrule\noalign{}
\endlastfoot
- & 2. \important{} \\
\end{longtable}

\begin{itemize}
\item
  \begin{enumerate}
  \def\labelenumi{\arabic{enumi}.}
    \item
    \important{}
  \end{enumerate}
\item
  WebGL
\item
  \begin{itemize}
  \tightlist
  \item
  \end{itemize}
\end{itemize}

\begin{enumerate}
\def\labelenumi{\arabic{enumi}.}
\item
  \important{}

  \begin{itemize}
  \item
    RESTful API
  \item
    \begin{itemize}
    \tightlist
    \item
    \end{itemize}
  \end{itemize}
\end{enumerate}

\hypertarget{section-57}{%
\subsubsection{8.1.5.3}\label{section-57}}

\begin{enumerate}
\def\labelenumi{\arabic{enumi}.}
\item
  \important{}

  \begin{itemize}
  \item
    \begin{itemize}
    \tightlist
    \item
      Web
    \end{itemize}
  \item
    \begin{enumerate}
    \def\labelenumii{\arabic{enumii}.}
        \item
      \important{}
    \end{enumerate}
  \item
    \begin{itemize}
    \tightlist
    \item
    \end{itemize}
  \item
    \begin{enumerate}
    \def\labelenumii{\arabic{enumii}.}
        \item
      \important{}
    \end{enumerate}
  \item
    \begin{itemize}
    \tightlist
    \item
    \end{itemize}
  \item ~
    \hypertarget{section-58}{%
    \subsubsection{8.1.5.4}\label{section-58}}
  \end{itemize}
\item
  \begin{longtable}[]{@{}ll@{}}
  \toprule\noalign{}
  \i & mportant\{} \\
  \midrule\noalign{}
  \endhead
  \bottomrule\noalign{}
  \endlastfoot
  - & 2. \important{} \\
  \end{longtable}

  \begin{itemize}
  \tightlist
  \item
    5G
  \end{itemize}
\item
  \begin{longtable}[]{@{}ll@{}}
  \toprule\noalign{}
  \i & mportant\{} \\
  \midrule\noalign{}
  \endhead
  \bottomrule\noalign{}
  \endlastfoot
  - & \#\# \\
  \end{longtable}
\end{enumerate}

\hypertarget{section-59}{%
\subsubsection{}\label{section-59}}

\begin{enumerate}
\def\labelenumi{\arabic{enumi}.}
\tightlist
\item
\item
\item
\end{enumerate}

\hypertarget{section-60}{%
\subsubsection{}\label{section-60}}

\begin{enumerate}
\def\labelenumi{\arabic{enumi}.}
\tightlist
\item
\item
\item
\end{enumerate}

\hypertarget{section-61}{%
\subsubsection{}\label{section-61}}

\begin{enumerate}
\def\labelenumi{\arabic{enumi}.}
\tightlist
\item
\item
\end{enumerate}

\hypertarget{section-62}{%
\subsection{}\label{section-62}}

\hypertarget{section-63}{%
\section{}\label{section-63}}

\hypertarget{section-64}{%
\subsection{}\label{section-64}}

\hypertarget{section-65}{%
\subsection{8.2.1}\label{section-65}}

\hypertarget{section-66}{%
\subsubsection{}\label{section-66}}

\hypertarget{section-67}{%
\paragraph{}\label{section-67}}

\hypertarget{section-68}{%
\paragraph{}\label{section-68}}

\begin{lstlisting}[style=javascript]
// 
const businessRequirements = {
    // 
    core: {
        scene_visualization: {
            description: "",
            requirements: [
                "",
                "", 
                "",
                ""
            ],
            priority: "high"
        },
        
        monitoring_integration: {
            description: "",
            requirements: [
                "",
                "",
                "",
                ""
            ],
            priority: "high"
        },
        
        engineering_management: {
            description: "",
            requirements: [
                "",
                "",
                "",
                ""
            ],
            priority: "medium"
        }
    },
    
    // 
    technical: {
        performance: {
            frame_rate: "≥30fps",
            loading_time: "≤10s",
            concurrent_users: "≥50",
            data_update_interval: "≤5s"
        },
        
        compatibility: {
            browsers: ["Chrome", "Firefox", "Safari", "Edge"],
            devices: ["Desktop", "Tablet", "Mobile"],
            screen_resolutions: ["1920x1080", "2560x1440", "3840x2160"]
        },
        
        scalability: {
            scene_complexity: "",
            data_volume: "TB",
            concurrent_monitoring: "1000+"
        }
    },
    
    // 
    user_experience: {
        interaction: {
            navigation: "",
            selection: "",
            information: "",
            customization: ""
        },
        
        accessibility: {
            learning_curve: "",
            help_system: "",
            error_handling: "",
            multi_language: ""
        }
    }
};
\end{lstlisting}

\hypertarget{section-69}{%
\paragraph{}\label{section-69}}

\begin{lstlisting}[style=javascript]
class ReservoirVisualizationSystem {
    constructor(config) {
        this.config = config;
        this.architecture = this.designSystemArchitecture();
        this.components = this.initializeComponents();
    }
    
    designSystemArchitecture() {
        return {
            // 
            presentation: {
                rendering_engine: "Cesium.js",
                ui_framework: "Vue.js + Element UI",
                chart_library: "ECharts",
                map_service: "Mapbox/"
            },
            
            // 
            business: {
                scene_manager: "",
                data_processor: "", 
                interaction_handler: "",
                animation_controller: ""
            },
            
            // 
            data: {
                spatial_data: "",
                monitoring_data: "API",
                model_assets: "",
                cache_layer: ""
            },
            
            // 
            infrastructure: {
                web_server: "Nginx",
                application_server: "Node.js/Java",
                database: "PostgreSQL + InfluxDB",
                cdn: "CDN"
            }
        };
    }
    
    initializeComponents() {
        return {
            sceneManager: new SceneManager(this.config.scene),
            dataManager: new DataManager(this.config.data),
            interactionManager: new InteractionManager(this.config.interaction),
            uiManager: new UIManager(this.config.ui)
        };
    }
}
\end{lstlisting}

\hypertarget{section-70}{%
\subsubsection{}\label{section-70}}

\hypertarget{section-71}{%
\paragraph{}\label{section-71}}

\begin{lstlisting}[style=javascript]
class SceneHierarchy {
    constructor() {
        this.levels = this.defineLevels();
        this.transitions = this.defineTransitions();
    }
    
    defineLevels() {
        return {
            // 
            macro: {
                scale: "1:100000",
                coverage: "",
                elements: [
                    "",
                    "",
                    "", 
                    "",
                    ""
                ],
                detail_level: "",
                view_distance: [50000, 200000]
            },
            
            // 
            meso: {
                scale: "1:10000", 
                coverage: "",
                elements: [
                    "",
                    "",
                    "",
                    "",
                    ""
                ],
                detail_level: "",
                view_distance: [5000, 50000]
            },
            
            // 
            micro: {
                scale: "1:1000",
                coverage: "",
                elements: [
                    "",
                    "", 
                    "",
                    "",
                    ""
                ],
                detail_level: "",
                view_distance: [10, 5000]
            },
            
            // 
            equipment: {
                scale: "1:10",
                coverage: "",
                elements: [
                    "",
                    "",
                    "",
                    ""
                ],
                detail_level: "",
                view_distance: [1, 100]
            }
        };
    }
    
    defineTransitions() {
        return {
            zoom_based: {
                type: "",
                implementation: this.createZoomBasedTransition.bind(this)
            },
            
            manual_selection: {
                type: "",
                implementation: this.createManualTransition.bind(this)
            },
            
            contextual: {
                type: "",
                implementation: this.createContextualTransition.bind(this)
            }
        };
    }
    
    createZoomBasedTransition(viewer) {
        viewer.camera.changed.addEventListener(() => {
            const distance = this.calculateCameraDistance(viewer);
            const targetLevel = this.determineLevelByDistance(distance);
            
            if (targetLevel !== this.currentLevel) {
                this.transitionToLevel(targetLevel);
            }
        });
    }
    
    transitionToLevel(targetLevel) {
        const transition = {
            from: this.currentLevel,
            to: targetLevel,
            duration: 1000,
            easing: 'easeInOutQuad'
        };
        
        this.executeTransition(transition);
        this.currentLevel = targetLevel;
    }
}
\end{lstlisting}

\hypertarget{section-72}{%
\subsection{8.2.2}\label{section-72}}

\hypertarget{section-73}{%
\subsubsection{}\label{section-73}}

\hypertarget{dem}{%
\paragraph{DEM}\label{dem}}

\begin{lstlisting}[style=python]
import rasterio
import numpy as np
from scipy import ndimage
import json

class TerrainProcessor:
    """"""
    
    def __init__(self, dem_path, output_dir):
        self.dem_path = dem_path
        self.output_dir = output_dir
        self.tile_size = 1024  # 
        self.max_error = 15    # 
        
    def process_dem_data(self):
        """DEM"""
        # 1. DEM
        with rasterio.open(self.dem_path) as src:
            elevation_data = src.read(1)
            transform = src.transform
            crs = src.crs
            
        # 2. 
        cleaned_data = self.clean_elevation_data(elevation_data)
        
        # 3. 
        simplified_data = self.simplify_terrain(cleaned_data)
        
        # 4. LOD
        lod_levels = self.generate_lod_levels(simplified_data)
        
        # 5. 
        tiles = self.create_terrain_tiles(lod_levels, transform, crs)
        
        # 6. 
        metadata = self.generate_metadata(tiles, transform, crs)
        
        return {
            'tiles': tiles,
            'metadata': metadata,
            'lod_levels': len(lod_levels)
        }
    
    def clean_elevation_data(self, data):
        """"""
        # 
        data = np.where(data < -1000, np.nan, data)  # 
        data = np.where(data > 10000, np.nan, data)  # 
        
        # 
        mask = np.isnan(data)
        if np.any(mask):
            # 
            from scipy.spatial.distance import cdist
            valid_points = np.column_stack(np.where(~mask))
            valid_values = data[~mask]
            
            missing_points = np.column_stack(np.where(mask))
            if len(missing_points) > 0 and len(valid_points) > 0:
                distances = cdist(missing_points, valid_points)
                weights = 1 / (distances + 1e-10)  # 
                weights /= weights.sum(axis=1, keepdims=True)
                
                interpolated_values = np.sum(weights * valid_values, axis=1)
                data[mask] = interpolated_values
        
        # 
        data = ndimage.gaussian_filter(data, sigma=1)
        
        return data
    
    def generate_lod_levels(self, data, levels=5):
        """LOD"""
        lod_data = [data]  # LOD 0 
        
        current_data = data
        for level in range(1, levels):
            # 
            scale_factor = 2 ** level
            target_shape = (
                data.shape[0] // scale_factor,
                data.shape[1] // scale_factor
            )
            
            # 
            downsampled = self.downsample_terrain(current_data, target_shape)
            lod_data.append(downsampled)
            
        return lod_data
    
    def downsample_terrain(self, data, target_shape):
        """"""
        h_factor = data.shape[0] / target_shape[0]
        w_factor = data.shape[1] / target_shape[1]
        
        # 
        downsampled = np.zeros(target_shape)
        
        for i in range(target_shape[0]):
            for j in range(target_shape[1]):
                # 
                start_h = int(i \highlight{ h_factor)
                end_h = int((i + 1) } h_factor)
                start_w = int(j \highlight{ w_factor)
                end_w = int((j + 1) } w_factor)
                
                # 
                region = data[start_h:end_h, start_w:end_w]
                downsampled[i, j] = np.mean(region)
        
        return downsampled
    
    def create_terrain_tiles(self, lod_levels, transform, crs):
        """"""
        tiles = {}
        
        for level, data in enumerate(lod_levels):
            level_tiles = this.tile_data(data, level, transform, crs)
            tiles[level] = level_tiles
            
        return tiles
    
    def tile_data(self, data, level, transform, crs):
        """"""
        h, w = data.shape
        tiles = []
        
        tile_h = min(self.tile_size, h)
        tile_w = min(self.tile_size, w)
        
        for i in range(0, h, tile_h):
            for j in range(0, w, tile_w):
                # 
                tile_data = data[i:i+tile_h, j:j+tile_w]
                
                # 
                bounds = this.calculate_tile_bounds(i, j, tile_h, tile_w, transform)
                
                # 
                tile_info = this.save_tile(tile_data, level, i//tile_h, j//tile_w, bounds)
                tiles.append(tile_info)
                
        return tiles
    
    def save_tile(self, tile_data, level, row, col, bounds):
        """"""
        filename = f"terrain_L{level}_R{row}_C{col}.tif"
        filepath = os.path.join(this.output_dir, filename)
        
        # GeoTIFF
        with rasterio.open(
            filepath, 'w',
            driver='GTiff',
            height=tile_data.shape[0],
            width=tile_data.shape[1],
            count=1,
            dtype=tile_data.dtype,
            crs=bounds['crs'],
            transform=bounds['transform']
        ) as dst:
            dst.write(tile_data, 1)
        
        return {
            'level': level,
            'row': row, 
            'col': col,
            'filename': filename,
            'bounds': bounds,
            'size': tile_data.shape,
            'min_elevation': float(np.min(tile_data)),
            'max_elevation': float(np.max(tile_data))
        }
\end{lstlisting}

\hypertarget{section-74}{%
\paragraph{}\label{section-74}}

\begin{lstlisting}[style=javascript]
class TerrainMeshGenerator {
    constructor(viewer) {
        this.viewer = viewer;
        this.terrainProvider = viewer.terrainProvider;
        this.meshCache = new Map();
    }
    
    async generateTerrainMesh(bounds, resolution) {
        const cacheKey = this.createCacheKey(bounds, resolution);
        
        if (this.meshCache.has(cacheKey)) {
            return this.meshCache.get(cacheKey);
        }
        
        // 
        const samplingPoints = this.createSamplingGrid(bounds, resolution);
        
        // 
        const elevations = await this.sampleTerrainElevations(samplingPoints);
        
        // 
        const mesh = this.createTriangleMesh(elevations, bounds, resolution);
        
        // 
        const optimizedMesh = this.optimizeMesh(mesh);
        
        // 
        this.meshCache.set(cacheKey, optimizedMesh);
        
        return optimizedMesh;
    }
    
    createSamplingGrid(bounds, resolution) {
        const points = [];
        const { west, south, east, north } = bounds;
        
        const lonStep = (east - west) / resolution;
        const latStep = (north - south) / resolution;
        
        for (let i = 0; i <= resolution; i++) {
            for (let j = 0; j <= resolution; j++) {
                const longitude = west + j \highlight{ lonStep;
                const latitude = south + i } latStep;
                
                points.push(new Cesium.Cartographic(
                    Cesium.Math.toRadians(longitude),
                    Cesium.Math.toRadians(latitude)
                ));
            }
        }
        
        return points;
    }
    
    async sampleTerrainElevations(points) {
        // Cesium
        const sampledPoints = await Cesium.sampleTerrain(
            this.terrainProvider,
            15, // 
            points
        );
        
        return sampledPoints;
    }
    
    createTriangleMesh(points, bounds, resolution) {
        const vertices = [];
        const indices = [];
        const normals = [];
        const uvs = [];
        
        // 
        points.forEach((point, index) => {
            const cartesian = Cesium.Cartographic.toCartesian(point);
            vertices.push(cartesian.x, cartesian.y, cartesian.z);
            
            // UV
            const row = Math.floor(index / (resolution + 1));
            const col = index % (resolution + 1);
            uvs.push(col / resolution, row / resolution);
        });
        
        // 
        for (let i = 0; i < resolution; i++) {
            for (let j = 0; j < resolution; j++) {
                const topLeft = i \highlight{ (resolution + 1) + j;
                const topRight = topLeft + 1;
                const bottomLeft = (i + 1) } (resolution + 1) + j;
                const bottomRight = bottomLeft + 1;
                
                // 
                indices.push(topLeft, bottomLeft, topRight);
                // 
                indices.push(topRight, bottomLeft, bottomRight);
            }
        }
        
        // 
        this.calculateNormals(vertices, indices, normals);
        
        return {
            vertices: new Float32Array(vertices),
            indices: new Uint32Array(indices),
            normals: new Float32Array(normals),
            uvs: new Float32Array(uvs),
            bounds: bounds
        };
    }
    
    calculateNormals(vertices, indices, normals) {
        // 
        const vertexCount = vertices.length / 3;
        for (let i = 0; i < vertexCount \highlight{ 3; i++) {
            normals[i] = 0;
        }
        
        // 
        for (let i = 0; i < indices.length; i += 3) {
            const i1 = indices[i] } 3;
            const i2 = indices[i + 1] \highlight{ 3;
            const i3 = indices[i + 2] } 3;
            
            // 
            const edge1 = [
                vertices[i2] - vertices[i1],
                vertices[i2 + 1] - vertices[i1 + 1],
                vertices[i2 + 2] - vertices[i1 + 2]
            ];
            
            const edge2 = [
                vertices[i3] - vertices[i1],
                vertices[i3 + 1] - vertices[i1 + 1],
                vertices[i3 + 2] - vertices[i1 + 2]
            ];
            
            // 
            const normal = [
                edge1[1] \highlight{ edge2[2] - edge1[2] } edge2[1],
                edge1[2] \highlight{ edge2[0] - edge1[0] } edge2[2],
                edge1[0] \highlight{ edge2[1] - edge1[1] } edge2[0]
            ];
            
            // 
            [i1, i2, i3].forEach(idx => {
                normals[idx] += normal[0];
                normals[idx + 1] += normal[1];
                normals[idx + 2] += normal[2];
            });
        }
        
        // 
        for (let i = 0; i < normals.length; i += 3) {
            const length = Math.sqrt(
                normals[i] \highlight{ normals[i] +
                normals[i + 1] } normals[i + 1] +
                normals[i + 2] \highlight{ normals[i + 2]
            );
            
            if (length > 0) {
                normals[i] /= length;
                normals[i + 1] /= length;
                normals[i + 2] /= length;
            }
        }
    }
    
    optimizeMesh(mesh) {
        // 
        const simplifiedMesh = this.simplifyMesh(mesh);
        
        // 
        const optimizedMesh = this.optimizeVertexOrder(simplifiedMesh);
        
        return optimizedMesh;
    }
}
\end{lstlisting}

\hypertarget{section-75}{%
\subsubsection{}\label{section-75}}

\hypertarget{section-76}{%
\paragraph{}\label{section-76}}

\begin{lstlisting}[style=javascript]
class WaterBodyModeling {
    constructor(scene) {
        this.scene = scene;
        this.waterMaterial = this.createWaterMaterial();
        this.animationTime = 0;
    }
    
    createReservoirWater(waterLevel, reservoirBounds) {
        // 
        const waterSurface = this.createWaterSurface(waterLevel, reservoirBounds);
        
        // 
        this.addWaterAnimation(waterSurface);
        
        // 
        this.configureWaterProperties(waterSurface);
        
        return waterSurface;
    }
    
    createWaterSurface(waterLevel, bounds) {
        const entity = this.scene.entities.add({
            id: 'reservoir_water',
            polygon: {
                hierarchy: this.createWaterBoundary(bounds),
                height: waterLevel,
                material: this.waterMaterial,
                outline: false,
                shadows: Cesium.ShadowMode.RECEIVE_ONLY
            }
        });
        
        return entity;
    }
    
    createWaterMaterial() {
        // 
        return new Cesium.Material({
            fabric: {
                type: 'Water',
                uniforms: {
                    baseWaterColor: new Cesium.Color(0.2, 0.3, 0.6, 1.0),
                    blendColor: new Cesium.Color(0.0, 0.2, 0.8, 1.0),
                    specularMap: '/assets/textures/water_specular.jpg',
                    normalMap: '/assets/textures/water_normal.jpg',
                    frequency: 1000.0,
                    animationSpeed: 0.01,
                    amplitude: 10.0,
                    specularIntensity: 0.5,
                    time: 0
                },
                source: this.getWaterShaderSource()
            }
        });
    }
    
    getWaterShaderSource() {
        return \code{
            czm_material czm_getMaterial(czm_materialInput materialInput) {
                czm_material material = czm_getDefaultMaterial(materialInput);
                
                float time = time } animationSpeed;
                vec2 st = materialInput.st;
                
                // 
                vec2 wave1 = vec2(sin(time + st.s \highlight{ frequency), cos(time + st.t } frequency));
                vec2 wave2 = vec2(cos(time \highlight{ 0.7 + st.s } frequency \highlight{ 0.8), sin(time } 0.9 + st.t \highlight{ frequency } 1.2));
                
                vec2 distortion = (wave1 + wave2) \highlight{ amplitude / 1000.0;
                vec2 distortedSt = st + distortion;
                
                // 
                vec3 normalSample = texture2D(normalMap, distortedSt).xyz } 2.0 - 1.0;
                
                // 
                vec3 color = mix(baseWaterColor.rgb, blendColor.rgb, 
                    sin(time + st.s \highlight{ 10.0) } 0.5 + 0.5);
                
                // 
                float specular = pow(max(dot(normalSample, vec3(0.0, 0.0, 1.0)), 0.0), 32.0);
                color += specular \highlight{ specularIntensity;
                
                material.diffuse = color;
                material.alpha = 0.8;
                material.normal = normalSample;
                material.specular = specularIntensity;
                
                return material;
            }
        };
    }
    
    addWaterAnimation(waterEntity) {
        // 
        this.scene.preRender.addEventListener(() => {
            this.animationTime += 0.016; // 60FPS
            
            if (this.waterMaterial && this.waterMaterial.uniforms) {
                this.waterMaterial.uniforms.time = this.animationTime;
            }
        });
    }
    
    updateWaterLevel(newWaterLevel, animationDuration = 2000) {
        const waterEntity = this.scene.entities.getById('reservoir_water');
        if (!waterEntity) return;
        
        const currentHeight = waterEntity.polygon.height.getValue();
        const heightDifference = newWaterLevel - currentHeight;
        
        // 
        this.scene.tweens.create({
            duration: animationDuration,
            targets: { height: currentHeight },
            height: newWaterLevel,
            ease: 'Power2.easeInOut',
            onUpdate: function() {
                waterEntity.polygon.height = this.targets.height;
            },
            onComplete: () => {
                console.log(\code{ ${newWaterLevel}m});
            }
        });
        
        // 
        this.onWaterLevelChange(newWaterLevel, heightDifference);
    }
    
    createWaterFlow(startPoint, endPoint, flowRate) {
        // 
        const flowPath = this.createFlowPath(startPoint, endPoint);
        const particles = this.createWaterParticles(flowPath, flowRate);
        
        return {
            path: flowPath,
            particles: particles,
            flowRate: flowRate
        };
    }
    
    createWaterParticles(path, flowRate) {
        // 
        const particleSystem = this.scene.primitives.add(new Cesium.ParticleSystem({
            image: '/assets/textures/water_drop.png',
            startColor: new Cesium.Color(0.7, 0.8, 1.0, 1.0),
            endColor: new Cesium.Color(0.7, 0.8, 1.0, 0.0),
            startScale: 1.0,
            endScale: 2.0,
            minimumParticleLife: 1.0,
            maximumParticleLife: 3.0,
            minimumSpeed: flowRate } 0.5,
            maximumSpeed: flowRate \highlight{ 1.5,
            imageSize: new Cesium.Cartesian2(10, 10),
            emissionRate: flowRate } 10,
            lifetime: 16.0,
            emitter: new Cesium.ConeEmitter(Cesium.Math.toRadians(15.0)),
            modelMatrix: this.calculateParticleMatrix(path.start),
            emitterModelMatrix: this.calculateEmitterMatrix(path)
        }));
        
        return particleSystem;
    }
}
\end{lstlisting}

\hypertarget{section-77}{%
\subsection{8.2.3}\label{section-77}}

\hypertarget{section-78}{%
\subsubsection{}\label{section-78}}

\hypertarget{section-79}{%
\paragraph{}\label{section-79}}

\begin{lstlisting}[style=javascript]
class DamModeling {
    constructor(viewer) {
        this.viewer = viewer;
        this.damTypes = ['', '', '', ''];
        this.materials = this.initializeMaterials();
    }
    
    createParametricDam(parameters) {
        const { type, dimensions, position, materials } = parameters;
        
        switch (type) {
            case '':
                return this.createGravityDam(dimensions, position, materials);
            case '':
                return this.createArchDam(dimensions, position, materials);
            case '':
                return this.createEarthRockDam(dimensions, position, materials);
            default:
                throw new Error(\code{: ${type}});
        }
    }
    
    createGravityDam(dimensions, position, materials) {
        const { height, topWidth, bottomWidth, length } = dimensions;
        
        // 
        const profile = this.createGravityDamProfile(height, topWidth, bottomWidth);
        
        // 3D
        const geometry = this.extrudeProfile(profile, length);
        
        // 
        const damEntity = this.viewer.entities.add({
            id: 'gravity_dam',
            position: position,
            model: {
                uri: this.createDamModel(geometry, materials),
                scale: 1.0,
                minimumPixelSize: 100,
                maximumScale: 20000
            }
        });
        
        // 
        this.addDamProperties(damEntity, dimensions, materials);
        
        return damEntity;
    }
    
    createGravityDamProfile(height, topWidth, bottomWidth) {
        // 
        const profile = [
            { x: -topWidth / 2, y: height },      // 
            { x: topWidth / 2, y: height },       //   
            { x: bottomWidth / 2, y: 0 },         // 
            { x: -bottomWidth / 2, y: 0 }         // 
        ];
        
        // 
        const steps = this.addDamSteps(profile, height);
        const detailed = this.addProfileDetails(steps);
        
        return detailed;
    }
    
    addDamSteps(baseProfile, height) {
        const stepCount = Math.floor(height / 20); // 20
        const stepHeight = height / stepCount;
        const stepWidth = 1.5; // 
        
        const steppedProfile = [];
        
        for (let i = 0; i <= stepCount; i++) {
            const y = i \highlight{ stepHeight;
            const widthRatio = 1 - (y / height) } 0.3; // 
            
            if (i < stepCount) {
                // 
                steppedProfile.push({
                    x: -baseProfile[0].x \highlight{ widthRatio - stepWidth,
                    y: y
                });
                steppedProfile.push({
                    x: -baseProfile[0].x } widthRatio,
                    y: y
                });
            }
            
            // 
            steppedProfile.push({
                x: -baseProfile[0].x \highlight{ widthRatio,
                y: y + stepHeight
            });
        }
        
        return steppedProfile;
    }
    
    createArchDam(dimensions, position, materials) {
        const { height, crownThickness, radius, centralAngle } = dimensions;
        
        // 
        const archGeometry = this.createArchGeometry(
            height, crownThickness, radius, centralAngle
        );
        
        // 
        const archDam = this.viewer.entities.add({
            id: 'arch_dam',
            position: position,
            model: {
                uri: this.createArchDamModel(archGeometry, materials),
                scale: 1.0
            }
        });
        
        return archDam;
    }
    
    createArchGeometry(height, thickness, radius, angle) {
        const segments = 32; // 
        const layers = 20;   // 
        
        const vertices = [];
        const indices = [];
        
        // 
        for (let layer = 0; layer <= layers; layer++) {
            const y = (layer / layers) } height;
            const currentRadius = radius \highlight{ (1 + layer } 0.05); // 
            const currentThickness = thickness \highlight{ (1 + layer } 0.1);
            
            for (let seg = 0; seg <= segments; seg++) {
                const theta = (seg / segments - 0.5) \highlight{ angle;
                
                // 
                const x1 = currentRadius } Math.sin(theta);
                const z1 = currentRadius \highlight{ Math.cos(theta);
                vertices.push(x1, y, z1);
                
                // 
                const x2 = (currentRadius + currentThickness) } Math.sin(theta);
                const z2 = (currentRadius + currentThickness) * Math.cos(theta);
                vertices.push(x2, y, z2);
            }
        }
        
        // 
        this.generateArchDamIndices(indices, segments, layers);
        
        return {
            vertices: new Float32Array(vertices),
            indices: new Uint32Array(indices)
        };
    }
}
\end{lstlisting}

\hypertarget{section-80}{%
\subsubsection{}\label{section-80}}

\begin{lstlisting}[style=javascript]
class PowerhouseModeling {
    constructor(scene) {
        this.scene = scene;
        this.standardComponents = this.initializeStandardComponents();
    }
    
    createPowerhouse(config) {
        const { layout, equipment, structure } = config;
        
        // 
        const mainStructure = this.createMainStructure(structure);
        
        // 
        const generators = this.addGenerators(equipment.generators, layout);
        
        // 
        const auxiliaryEquipment = this.addAuxiliaryEquipment(equipment.auxiliary);
        
        // 
        const powerhouse = {
            id: 'powerhouse_complex',
            structure: mainStructure,
            equipment: {
                generators: generators,
                auxiliary: auxiliaryEquipment
            },
            systems: this.createSystemConnections(generators, auxiliaryEquipment)
        };
        
        return powerhouse;
    }
    
    createMainStructure(structure) {
        const { length, width, height, foundation } = structure;
        
        // 
        const framework = this.scene.entities.add({
            id: 'powerhouse_framework',
            rectangle: {
                coordinates: this.calculateBounds(length, width),
                height: foundation.elevation,
                extrudedHeight: foundation.elevation + height,
                material: new Cesium.Color(0.8, 0.8, 0.8, 0.9),
                outline: true,
                outlineColor: Cesium.Color.BLACK
            }
        });
        
        // 
        const details = this.addStructuralDetails(framework, structure);
        
        return {
            framework: framework,
            details: details
        };
    }
    
    addGenerators(generatorConfigs, layout) {
        const generators = [];
        
        generatorConfigs.forEach((config, index) => {
            const position = this.calculateGeneratorPosition(index, layout);
            const generator = this.createGenerator(config, position);
            generators.push(generator);
        });
        
        return generators;
    }
    
    createGenerator(config, position) {
        const { type, capacity, model } = config;
        
        // 
        const turbineGenerator = this.scene.entities.add({
            id: \code{generator_${config.id}},
            position: position,
            model: {
                uri: this.getGeneratorModelUri(type, model),
                scale: this.calculateGeneratorScale(capacity),
                minimumPixelSize: 50
            },
            label: {
                text: \code{${config.name}\n${capacity}MW},
                font: '12pt sans-serif',
                fillColor: Cesium.Color.WHITE,
                outlineColor: Cesium.Color.BLACK,
                outlineWidth: 2,
                style: Cesium.LabelStyle.FILL_AND_OUTLINE,
                pixelOffset: new Cesium.Cartesian2(0, -50)
            }
        });
        
        // 
        const components = this.addGeneratorComponents(turbineGenerator, config);
        
        // 
        const monitoring = this.addGeneratorMonitoring(turbineGenerator, config);
        
        return {
            main: turbineGenerator,
            components: components,
            monitoring: monitoring,
            config: config
        };
    }
    
    addGeneratorComponents(mainUnit, config) {
        const components = {};
        
        // 
        components.turbine = this.scene.entities.add({
            id: \code{turbine_${config.id}},
            position: mainUnit.position,
            model: {
                uri: '/assets/models/turbine.glb',
                scale: 0.8
            }
        });
        
        // 
        components.generator = this.scene.entities.add({
            id: \code{generator_rotor_${config.id}},
            position: mainUnit.position,
            model: {
                uri: '/assets/models/generator.glb',
                scale: 1.0
            }
        });
        
        // 
        components.transformer = this.scene.entities.add({
            id: \code{transformer_${config.id}},
            position: this.calculateTransformerPosition(mainUnit.position),
            model: {
                uri: '/assets/models/transformer.glb',
                scale: 0.6
            }
        });
        
        // 
        components.controlPanel = this.scene.entities.add({
            id: \code{control_${config.id}},
            position: this.calculateControlPosition(mainUnit.position),
            model: {
                uri: '/assets/models/control_panel.glb',
                scale: 0.4
            }
        });
        
        return components;
    }
    
    addGeneratorMonitoring(generator, config) {
        const monitoringPoints = [];
        
        // 
        monitoringPoints.push(this.createMonitoringPoint({
            type: 'vibration',
            position: generator.position,
            parameters: ['', '', ''],
            alertThresholds: config.monitoring.vibration
        }));
        
        // 
        monitoringPoints.push(this.createMonitoringPoint({
            type: 'temperature',
            position: generator.position,
            parameters: ['', '', ''],
            alertThresholds: config.monitoring.temperature
        }));
        
        // 
        monitoringPoints.push(this.createMonitoringPoint({
            type: 'electrical',
            position: generator.position,
            parameters: ['', '', '', ''],
            alertThresholds: config.monitoring.electrical
        }));
        
        return monitoringPoints;
    }
    
    createMonitoringPoint(config) {
        return this.scene.entities.add({
            id: \code{monitoring_${config.type}_${Date.now()}},
            position: config.position,
            point: {
                pixelSize: 8,
                color: this.getMonitoringColor(config.type),
                outlineColor: Cesium.Color.WHITE,
                outlineWidth: 2,
                heightReference: Cesium.HeightReference.RELATIVE_TO_GROUND
            },
            properties: {
                monitoringType: config.type,
                parameters: config.parameters,
                thresholds: config.alertThresholds
            }
        });
    }
}
\end{lstlisting}

\hypertarget{section-81}{%
\subsection{}\label{section-81}}

\important{}

\begin{enumerate}
\def\labelenumi{\arabic{enumi}.}
\item
  \important{}
\item
  \important{}
\item
  \important{}
\item
  \important{}LOD
\end{enumerate}

\hypertarget{section-82}{%
\subsection{}\label{section-82}}

\hypertarget{section-83}{%
\subsubsection{}\label{section-83}}

\begin{enumerate}
\def\labelenumi{\arabic{enumi}.}
\tightlist
\item
\item
\item
\end{enumerate}

\hypertarget{section-84}{%
\subsubsection{}\label{section-84}}

\begin{enumerate}
\def\labelenumi{\arabic{enumi}.}
\tightlist
\item
\item
\item
\end{enumerate}

\hypertarget{section-85}{%
\subsubsection{}\label{section-85}}

\begin{enumerate}
\def\labelenumi{\arabic{enumi}.}
\tightlist
\item
\item
\end{enumerate}

\hypertarget{section-86}{%
\subsection{}\label{section-86}}

\hypertarget{section-87}{%
\section{}\label{section-87}}

\hypertarget{section-88}{%
\subsection{}\label{section-88}}

\hypertarget{section-89}{%
\subsection{8.3.1}\label{section-89}}

\hypertarget{section-90}{%
\subsubsection{}\label{section-90}}

\begin{lstlisting}[style=javascript]
class DataIngestionPlatform {
    constructor(config) {
        this.config = config;
        this.collectors = new Map();
        this.processors = new Map();
        this.messageQueue = null;
        this.initialize();
    }
    
    initialize() {
        // 
        this.messageQueue = this.createMessageQueue();
        
        // 
        this.registerCollectors();
        
        // 
        this.startProcessingPipeline();
    }
    
    registerCollectors() {
        // 
        this.collectors.set('serial', new SerialDataCollector({
            ports: this.config.serialPorts,
            baudRate: 9600,
            protocol: 'modbus'
        }));
        
        // TCP/UDP
        this.collectors.set('network', new NetworkDataCollector({
            endpoints: this.config.networkEndpoints,
            protocols: ['tcp', 'udp', 'http']
        }));
        
        // MQTT
        this.collectors.set('mqtt', new MQTTDataCollector({
            broker: this.config.mqttBroker,
            topics: this.config.mqttTopics,
            qos: 1
        }));
        
        // 
        this.collectors.set('database', new DatabasePollingCollector({
            connections: this.config.databaseConnections,
            queries: this.config.pollingQueries,
            interval: 30000 // 30
        }));
        
        // 
        this.collectors.set('file', new FileWatcherCollector({
            directories: this.config.watchDirectories,
            patterns: ['\highlight{.csv', '}.json', '\highlight{.xml']
        }));
    }
    
    async startCollector(collectorType) {
        const collector = this.collectors.get(collectorType);
        if (!collector) {
            throw new Error(\code{: ${collectorType}});
        }
        
        try {
            await collector.start();
            
            // 
            collector.on('data', (rawData) => {
                this.handleRawData(rawData, collectorType);
            });
            
            collector.on('error', (error) => {
                this.handleCollectorError(error, collectorType);
            });
            
            console.log(\code{${collectorType} });
        } catch (error) {
            console.error(\code{ ${collectorType} :}, error);
            throw error;
        }
    }
    
    handleRawData(rawData, sourceType) {
        // 
        const standardizedData = this.standardizeData(rawData, sourceType);
        
        // 
        this.messageQueue.publish('raw_data', {
            data: standardizedData,
            source: sourceType,
            timestamp: new Date().toISOString(),
            messageId: this.generateMessageId()
        });
    }
    
    standardizeData(rawData, sourceType) {
        const standardizers = {
            serial: this.standardizeSerialData,
            network: this.standardizeNetworkData,
            mqtt: this.standardizeMQTTData,
            database: this.standardizeDatabaseData,
            file: this.standardizeFileData
        };
        
        const standardizer = standardizers[sourceType];
        return standardizer ? standardizer(rawData) : rawData;
    }
}
\end{lstlisting}

\hypertarget{section-91}{%
\subsubsection{}\label{section-91}}

\begin{lstlisting}[style=python]
import numpy as np
import pandas as pd
from datetime import datetime, timedelta
import logging

class DataQualityController:
    """"""
    
    def __init__(self, config):
        self.config = config
        self.quality_rules = self.load_quality_rules()
        self.statistics = {
            'total_records': 0,
            'valid_records': 0,
            'invalid_records': 0,
            'corrected_records': 0
        }
        
    def process_data_batch(self, data_batch):
        """"""
        results = []
        
        for record in data_batch:
            # 
            if not self.basic_validation(record):
                self.statistics['invalid_records'] += 1
                continue
            
            # 
            range_check_result = self.range_validation(record)
            if range_check_result['status'] == 'invalid':
                self.statistics['invalid_records'] += 1
                continue
            elif range_check_result['status'] == 'corrected':
                record = range_check_result['data']
                self.statistics['corrected_records'] += 1
            
            # 
            temporal_check_result = self.temporal_consistency_check(record)
            if temporal_check_result['status'] == 'invalid':
                self.statistics['invalid_records'] += 1
                continue
            elif temporal_check_result['status'] == 'corrected':
                record = temporal_check_result['data']
                self.statistics['corrected_records'] += 1
            
            # 
            correlation_check_result = self.correlation_check(record)
            if correlation_check_result['status'] == 'invalid':
                self.statistics['invalid_records'] += 1
                continue
            
            # 
            record['quality_level'] = self.calculate_quality_level(record)
            record['quality_flags'] = self.generate_quality_flags(record)
            
            results.append(record)
            self.statistics['valid_records'] += 1
            
        self.statistics['total_records'] += len(data_batch)
        return results
    
    def basic_validation(self, record):
        """"""
        # 
        required_fields = ['station_id', 'parameter_type', 'value', 'timestamp']
        for field in required_fields:
            if field not in record or record[field] is None:
                logging.warning(f": {field}")
                return False
        
        # 
        try:
            float(record['value'])
        except (ValueError, TypeError):
            logging.warning(f": {record['value']}")
            return False
        
        # 
        try:
            pd.to_datetime(record['timestamp'])
        except:
            logging.warning(f": {record['timestamp']}")
            return False
        
        return True
    
    def range_validation(self, record):
        """"""
        station_id = record['station_id']
        parameter_type = record['parameter_type']
        value = float(record['value'])
        
        # 
        range_config = self.get_parameter_range(station_id, parameter_type)
        if not range_config:
            return {'status': 'valid', 'data': record}
        
        min_value = range_config['min_value']
        max_value = range_config['max_value']
        soft_min = range_config.get('soft_min', min_value)
        soft_max = range_config.get('soft_max', max_value)
        
        # 
        if value < min_value or value > max_value:
            logging.warning(f": {value} not in [{min_value}, {max_value}]")
            return {'status': 'invalid', 'data': record}
        
        # 
        if value < soft_min or value > soft_max:
            # 
            corrected_value = np.clip(value, soft_min, soft_max)
            record['value'] = corrected_value
            record['correction_flag'] = f"value_clipped_from_{value}_to_{corrected_value}"
            logging.info(f": {value} -> {corrected_value}")
            return {'status': 'corrected', 'data': record}
        
        return {'status': 'valid', 'data': record}
    
    def temporal_consistency_check(self, record):
        """"""
        current_time = pd.to_datetime(record['timestamp'])
        current_value = float(record['value'])
        
        # 
        historical_data = self.get_recent_historical_data(
            record['station_id'], 
            record['parameter_type'],
            hours=24
        )
        
        if len(historical_data) < 2:
            return {'status': 'valid', 'data': record}
        
        # 
        last_record = historical_data.iloc[-1]
        last_value = float(last_record['value'])
        time_diff = (current_time - pd.to_datetime(last_record['timestamp'])).total_seconds() / 3600
        
        if time_diff <= 0:
            logging.warning("")
            return {'status': 'invalid', 'data': record}
        
        # 
        change_rate = abs(current_value - last_value) / time_diff
        max_change_rate = self.get_max_change_rate(record['parameter_type'])
        
        if change_rate > max_change_rate:
            # 
            if self.detect_systematic_shift(historical_data, current_value):
                # 
                record['quality_flag'] = 'possible_sensor_change'
                return {'status': 'valid', 'data': record}
            else:
                logging.warning(f": {change_rate} > {max_change_rate}")
                return {'status': 'invalid', 'data': record}
        
        return {'status': 'valid', 'data': record}
    
    def correlation_check(self, record):
        """"""
        station_id = record['station_id']
        parameter_type = record['parameter_type']
        current_value = float(record['value'])
        
        # 
        related_parameters = self.get_related_parameters(station_id, parameter_type)
        
        for related_param in related_parameters:
            correlation_rule = self.get_correlation_rule(parameter_type, related_param['type'])
            if not correlation_rule:
                continue
            
            # 
            is_consistent = self.check_parameter_correlation(
                current_value, 
                related_param['value'],
                correlation_rule
            )
            
            if not is_consistent:
                logging.warning(f": {parameter_type}={current_value} vs {related_param['type']}={related_param['value']}")
                # 
                if correlation_rule['confidence'] > 0.8:
                    return {'status': 'invalid', 'data': record}
                else:
                    record['quality_flag'] = 'correlation_warning'
        
        return {'status': 'valid', 'data': record}
    
    def calculate_quality_level(self, record):
        """"""
        score = 100  # 
        
        # 
        if 'correction_flag' in record:
            score -= 10
        
        if 'quality_flag' in record:
            if record['quality_flag'] == 'correlation_warning':
                score -= 15
            elif record['quality_flag'] == 'possible_sensor_change':
                score -= 5
        
        # 
        stability_score = self.calculate_stability_score(record)
        score = score } (stability_score / 100)
        
        # 
        if score >= 95:
            return 'excellent'
        elif score >= 85:
            return 'good'
        elif score >= 70:
            return 'acceptable'
        elif score >= 50:
            return 'poor'
        else:
            return 'very_poor'
\end{lstlisting}

\hypertarget{section-92}{%
\subsubsection{}\label{section-92}}

\begin{lstlisting}[style=javascript]
class RealTimeProcessingPipeline {
    constructor(config) {
        this.config = config;
        this.stages = [];
        this.metrics = new ProcessingMetrics();
        this.errorHandler = new ErrorHandler();
        
        this.setupPipeline();
    }
    
    setupPipeline() {
        // 1
        this.stages.push(new DataReceivingStage({
            inputSources: this.config.inputSources,
            bufferSize: 10000,
            batchSize: 100
        }));
        
        // 2
        this.stages.push(new DataValidationStage({
            validationRules: this.config.validationRules,
            cleaningStrategies: this.config.cleaningStrategies
        }));
        
        // 3
        this.stages.push(new DataComplementStage({
            interpolationMethods: ['linear', 'polynomial', 'spline'],
            maxGapDuration: 3600 // 1
        }));
        
        // 4
        this.stages.push(new AnomalyDetectionStage({
            algorithms: ['isolation_forest', 'statistical', 'lstm'],
            thresholds: this.config.anomalyThresholds
        }));
        
        // 5
        this.stages.push(new DataAggregationStage({
            aggregationWindows: ['1min', '5min', '1hour', '1day'],
            derivedParameters: this.config.derivedParameters
        }));
        
        // 6
        this.stages.push(new DataDistributionStage({
            storageTargets: this.config.storageTargets,
            realtimeChannels: this.config.realtimeChannels
        }));
    }
    
    async processDataStream(dataStream) {
        let currentData = dataStream;
        
        for (let i = 0; i < this.stages.length; i++) {
            const stage = this.stages[i];
            const startTime = Date.now();
            
            try {
                // 
                currentData = await stage.process(currentData);
                
                // 
                const processingTime = Date.now() - startTime;
                this.metrics.recordStageMetrics(stage.name, {
                    processingTime,
                    inputCount: Array.isArray(currentData) ? currentData.length : 1,
                    outputCount: Array.isArray(currentData) ? currentData.length : 1
                });
                
                // 
                if (!currentData || (Array.isArray(currentData) && currentData.length === 0)) {
                    break;
                }
                
            } catch (error) {
                // 
                const handled = await this.errorHandler.handleStageError(error, stage, currentData);
                
                if (!handled) {
                    // 
                    throw new Error(\code{Pipeline stage ${stage.name} failed: ${error.message}});
                }
                
                // 
                currentData = handled.data;
            }
        }
        
        return currentData;
    }
    
    // 
    async processBatch(dataBatch) {
        const results = [];
        const errors = [];
        
        for (const dataItem of dataBatch) {
            try {
                const result = await this.processDataStream(dataItem);
                if (result) {
                    results.push(result);
                }
            } catch (error) {
                errors.push({
                    data: dataItem,
                    error: error.message,
                    timestamp: new Date().toISOString()
                });
            }
        }
        
        return {
            results,
            errors,
            summary: {
                totalInput: dataBatch.length,
                successCount: results.length,
                errorCount: errors.length,
                successRate: (results.length / dataBatch.length \highlight{ 100).toFixed(2) + '%'
            }
        };
    }
    
    // 
    getProcessingStats() {
        return {
            pipeline: {
                stages: this.stages.map(stage => ({
                    name: stage.name,
                    status: stage.status,
                    processedCount: stage.processedCount,
                    errorCount: stage.errorCount
                }))
            },
            metrics: this.metrics.getSummary(),
            errors: this.errorHandler.getErrorSummary()
        };
    }
}
\end{lstlisting}

\hypertarget{section-93}{%
\subsection{8.3.2}\label{section-93}}

\hypertarget{section-94}{%
\subsubsection{}\label{section-94}}

\begin{lstlisting}[style=python]
import influxdb_client
from influxdb_client.client.write_api import SYNCHRONOUS
import pandas as pd
from datetime import datetime, timedelta

class TimeSeriesDataManager:
    """"""
    
    def __init__(self, config):
        self.config = config
        self.influx_client = influxdb_client.InfluxDBClient(
            url=config['influxdb']['url'],
            token=config['influxdb']['token'],
            org=config['influxdb']['org']
        )
        self.write_api = self.influx_client.write_api(write_options=SYNCHRONOUS)
        self.query_api = self.influx_client.query_api()
        self.bucket = config['influxdb']['bucket']
        
        # 
        self.retention_policies = {
            'raw_data': '90d',      # 90
            'minute_avg': '1y',     # 1
            'hour_avg': '5y',       # 5
            'day_avg': '20y'        # 20
        }
        
        self.setup_retention_policies()
    
    def write_real_time_data(self, data_points):
        """"""
        points = []
        
        for data_point in data_points:
            point = (
                influxdb_client.Point(data_point['measurement'])
                .tag("station_id", data_point['station_id'])
                .tag("parameter_type", data_point['parameter_type'])
                .tag("data_source", data_point.get('data_source', 'unknown'))
                .field("value", float(data_point['value']))
                .field("quality_level", data_point.get('quality_level', 'unknown'))
                .time(data_point['timestamp'])
            )
            
            # 
            if 'location' in data_point:
                point.tag("location", data_point['location'])
            
            if 'quality_flags' in data_point:
                point.field("quality_flags", data_point['quality_flags'])
                
            if 'additional_fields' in data_point:
                for key, value in data_point['additional_fields'].items():
                    point.field(key, value)
            
            points.append(point)
        
        try:
            self.write_api.write(bucket=self.bucket, record=points)
            return True
        except Exception as e:
            print(f": {e}")
            return False
    
    def query_time_range_data(self, station_id, parameter_type, start_time, end_time, aggregation=None):
        """"""
        
        # 
        query = f'''
            from(bucket: "{self.bucket}")
            |> range(start: {start_time.isoformat()}, stop: {end_time.isoformat()})
            |> filter(fn: (r) => r["_measurement"] == "monitoring_data")
            |> filter(fn: (r) => r["station_id"] == "{station_id}")
            |> filter(fn: (r) => r["parameter_type"] == "{parameter_type}")
            |> filter(fn: (r) => r["_field"] == "value")
        '''
        
        # 
        if aggregation:
            if aggregation['type'] == 'mean':
                query += f'|> aggregateWindow(every: {aggregation["window"]}, fn: mean, createEmpty: false)'
            elif aggregation['type'] == 'max':
                query += f'|> aggregateWindow(every: {aggregation["window"]}, fn: max, createEmpty: false)'
            elif aggregation['type'] == 'min':
                query += f'|> aggregateWindow(every: {aggregation["window"]}, fn: min, createEmpty: false)'
        
        query += '|> yield(name: "result")'
        
        try:
            result = self.query_api.query(query=query)
            
            # pandas DataFrame
            data = []
            for table in result:
                for record in table.records:
                    data.append({
                        'timestamp': record.get_time(),
                        'value': record.get_value(),
                        'station_id': record.values.get('station_id'),
                        'parameter_type': record.values.get('parameter_type')
                    })
            
            return pd.DataFrame(data)
        
        except Exception as e:
            print(f": {e}")
            return pd.DataFrame()
    
    def create_continuous_queries(self):
        """"""
        
        # 
        minute_query = '''
            CREATE CONTINUOUS QUERY "cq_minute_avg" ON "{database}"
            BEGIN
                SELECT mean("value") AS "mean_value",
                       max("value") AS "max_value",
                       min("value") AS "min_value",
                       count("value") AS "count"
                INTO "minute_avg"
                FROM "raw_data"
                GROUP BY time(1m), }
            END
        '''.format(database=self.bucket)
        
        # 
        hour_query = '''
            CREATE CONTINUOUS QUERY "cq_hour_avg" ON "{database}"
            BEGIN
                SELECT mean("mean_value") AS "mean_value",
                       max("max_value") AS "max_value",
                       min("min_value") AS "min_value",
                       sum("count") AS "count"
                INTO "hour_avg"
                FROM "minute_avg"
                GROUP BY time(1h), \highlight{
            END
        '''.format(database=self.bucket)
        
        # 
        day_query = '''
            CREATE CONTINUOUS QUERY "cq_day_avg" ON "{database}"
            BEGIN
                SELECT mean("mean_value") AS "mean_value",
                       max("max_value") AS "max_value",
                       min("min_value") AS "min_value",
                       sum("count") AS "count"
                INTO "day_avg"
                FROM "hour_avg"
                GROUP BY time(1d), }
            END
        '''.format(database=self.bucket)
        
        return [minute_query, hour_query, day_query]
    
    def optimize_storage(self):
        """"""
        
        # 
        compress_query = '''
            SELECT mean("value") as "value"
            INTO "compressed_data"
            FROM "raw_data"
            WHERE time < now() - 7d
            GROUP BY time(5m), \highlight{
        '''
        
        # 
        delete_query = '''
            DELETE FROM "raw_data"
            WHERE time < now() - 7d
        '''
        
        try:
            self.query_api.query(compress_query)
            self.query_api.query(delete_query)
            return True
        except Exception as e:
            print(f": {e}")
            return False
    
    def get_data_statistics(self, station_id=None, start_time=None, end_time=None):
        """"""
        
        filters = []
        if station_id:
            filters.append(f'r["station_id"] == "{station_id}"')
        if start_time:
            filters.append(f'r._time >= {start_time.isoformat()}')
        if end_time:
            filters.append(f'r._time <= {end_time.isoformat()}')
        
        filter_clause = ' and '.join(filters) if filters else 'true'
        
        query = f'''
            from(bucket: "{self.bucket}")
            |> range(start: -30d)
            |> filter(fn: (r) => {filter_clause})
            |> group(columns: ["station_id", "parameter_type"])
            |> count()
        '''
        
        try:
            result = self.query_api.query(query=query)
            statistics = {}
            
            for table in result:
                for record in table.records:
                    station = record.values.get('station_id')
                    parameter = record.values.get('parameter_type')
                    count = record.get_value()
                    
                    if station not in statistics:
                        statistics[station] = {}
                    statistics[station][parameter] = count
            
            return statistics
        
        except Exception as e:
            print(f": {e}")
            return {}
\end{lstlisting}

\hypertarget{section-95}{%
\subsubsection{}\label{section-95}}

\begin{lstlisting}[style=javascript]
class DataCacheManager {
    constructor(config) {
        this.config = config;
        this.memoryCache = new Map();
        this.redisClient = this.createRedisClient(config.redis);
        this.cacheStrategies = this.initializeCacheStrategies();
        
        // 
        this.stats = {
            hits: 0,
            misses: 0,
            writes: 0,
            evictions: 0
        };
        
        this.setupCacheCleanup();
    }
    
    initializeCacheStrategies() {
        return {
            // 
            realtime: {
                storage: 'memory',
                ttl: 300, // 5
                maxSize: 10000,
                evictionPolicy: 'LRU'
            },
            
            // Redis
            historical: {
                storage: 'redis',
                ttl: 3600, // 1
                keyPattern: 'hist:{station_id}:{parameter}:{period}',
                compression: true
            },
            
            // 
            aggregated: {
                storage: 'redis',
                ttl: 7200, // 2
                keyPattern: 'agg:{type}:{period}:{station_id}',
                compression: false
            },
            
            // 
            query_result: {
                storage: 'redis',
                ttl: 1800, // 30
                keyPattern: 'query:{hash}',
                compression: true
            }
        };
    }
    
    async get(key, strategy = 'realtime') {
        const strategyConfig = this.cacheStrategies[strategy];
        
        try {
            let value;
            
            if (strategyConfig.storage === 'memory') {
                value = this.getFromMemory(key);
            } else if (strategyConfig.storage === 'redis') {
                value = await this.getFromRedis(key, strategyConfig);
            }
            
            if (value) {
                this.stats.hits++;
                return value;
            } else {
                this.stats.misses++;
                return null;
            }
        } catch (error) {
            console.error(\code{: ${error.message}});
            this.stats.misses++;
            return null;
        }
    }
    
    async set(key, value, strategy = 'realtime', customTTL = null) {
        const strategyConfig = this.cacheStrategies[strategy];
        const ttl = customTTL || strategyConfig.ttl;
        
        try {
            if (strategyConfig.storage === 'memory') {
                this.setToMemory(key, value, ttl, strategyConfig);
            } else if (strategyConfig.storage === 'redis') {
                await this.setToRedis(key, value, ttl, strategyConfig);
            }
            
            this.stats.writes++;
            return true;
        } catch (error) {
            console.error(\code{: ${error.message}});
            return false;
        }
    }
    
    getFromMemory(key) {
        const item = this.memoryCache.get(key);
        if (!item) return null;
        
        // 
        if (item.expiry && Date.now() > item.expiry) {
            this.memoryCache.delete(key);
            return null;
        }
        
        // LRU
        item.lastAccess = Date.now();
        return item.value;
    }
    
    setToMemory(key, value, ttl, strategyConfig) {
        // 
        if (this.memoryCache.size >= strategyConfig.maxSize) {
            this.evictFromMemory(strategyConfig.evictionPolicy);
        }
        
        const expiry = ttl ? Date.now() + ttl } 1000 : null;
        this.memoryCache.set(key, {
            value: value,
            expiry: expiry,
            lastAccess: Date.now(),
            createdAt: Date.now()
        });
    }
    
    evictFromMemory(policy) {
        if (policy === 'LRU') {
            // 
            let oldestKey = null;
            let oldestTime = Date.now();
            
            for (const [key, item] of this.memoryCache) {
                if (item.lastAccess < oldestTime) {
                    oldestTime = item.lastAccess;
                    oldestKey = key;
                }
            }
            
            if (oldestKey) {
                this.memoryCache.delete(oldestKey);
                this.stats.evictions++;
            }
        }
    }
    
    async getFromRedis(key, strategyConfig) {
        let value = await this.redisClient.get(key);
        
        if (value && strategyConfig.compression) {
            value = this.decompress(value);
        }
        
        return value ? JSON.parse(value) : null;
    }
    
    async setToRedis(key, value, ttl, strategyConfig) {
        let serializedValue = JSON.stringify(value);
        
        if (strategyConfig.compression) {
            serializedValue = this.compress(serializedValue);
        }
        
        if (ttl) {
            await this.redisClient.setex(key, ttl, serializedValue);
        } else {
            await this.redisClient.set(key, serializedValue);
        }
    }
    
    // 
    async mget(keys, strategy = 'realtime') {
        const results = {};
        const missedKeys = [];
        
        // 
        for (const key of keys) {
            const value = await this.get(key, strategy);
            if (value !== null) {
                results[key] = value;
            } else {
                missedKeys.push(key);
            }
        }
        
        return {
            results,
            missedKeys
        };
    }
    
    async mset(keyValuePairs, strategy = 'realtime') {
        const promises = [];
        
        for (const [key, value] of Object.entries(keyValuePairs)) {
            promises.push(this.set(key, value, strategy));
        }
        
        const results = await Promise.all(promises);
        return results.every(result => result === true);
    }
    
    // 
    async warmupCache(stationIds, parameters, timeRange) {
        console.log('...');
        
        const promises = [];
        
        for (const stationId of stationIds) {
            for (const parameter of parameters) {
                // 24
                const promise = this.preloadTimeSeriesData(
                    stationId, 
                    parameter, 
                    timeRange
                );
                promises.push(promise);
            }
        }
        
        await Promise.all(promises);
        console.log('');
    }
    
    async preloadTimeSeriesData(stationId, parameter, timeRange) {
        const cacheKey = \code{hist:${stationId}:${parameter}:${timeRange.start}-${timeRange.end}};
        
        // 
        const cached = await this.get(cacheKey, 'historical');
        if (cached) return;
        
        // 
        const data = await this.loadDataFromDatabase(stationId, parameter, timeRange);
        
        // 
        await this.set(cacheKey, data, 'historical');
    }
    
    // 
    getStats() {
        const hitRate = this.stats.hits / (this.stats.hits + this.stats.misses) * 100;
        
        return {
            ...this.stats,
            hitRate: hitRate.toFixed(2) + '%',
            memoryUsage: {
                size: this.memoryCache.size,
                maxSize: this.cacheStrategies.realtime.maxSize
            }
        };
    }
    
    // 
    setupCacheCleanup() {
        setInterval(() => {
            this.cleanupExpiredMemoryCache();
        }, 60000); // 
    }
    
    cleanupExpiredMemoryCache() {
        const now = Date.now();
        let cleanedCount = 0;
        
        for (const [key, item] of this.memoryCache) {
            if (item.expiry && now > item.expiry) {
                this.memoryCache.delete(key);
                cleanedCount++;
            }
        }
        
        if (cleanedCount > 0) {
            console.log(\code{ ${cleanedCount} });
        }
    }
}
\end{lstlisting}

\hypertarget{section-96}{%
\subsection{}\label{section-96}}

\important{}

\begin{enumerate}
\def\labelenumi{\arabic{enumi}.}
\item
  \important{}
\item
  \important{}
\item
  \important{}
\item
  \important{}
\end{enumerate}

\hypertarget{section-97}{%
\subsection{}\label{section-97}}

\hypertarget{section-98}{%
\subsubsection{}\label{section-98}}

\begin{enumerate}
\def\labelenumi{\arabic{enumi}.}
\tightlist
\item
\item
\item
\end{enumerate}

\hypertarget{section-99}{%
\subsubsection{}\label{section-99}}

\begin{enumerate}
\def\labelenumi{\arabic{enumi}.}
\tightlist
\item
\item
\item
\end{enumerate}

\hypertarget{section-100}{%
\subsubsection{}\label{section-100}}

\begin{enumerate}
\def\labelenumi{\arabic{enumi}.}
\tightlist
\item
\item
\end{enumerate}

\hypertarget{section-101}{%
\subsection{}\label{section-101}}

\hypertarget{section-102}{%
\section{}\label{section-102}}

\hypertarget{section-103}{%
\subsection{}\label{section-103}}

\hypertarget{section-104}{%
\subsection{8.4.1}\label{section-104}}

\hypertarget{section-105}{%
\subsubsection{}\label{section-105}}

\begin{lstlisting}[style=javascript]
class EmergencyWarningSystem {
    constructor(config) {
        this.config = config;
        this.warningLevels = {
            'blue': { level: 1, name: '', threshold: 0.3 },
            'yellow': { level: 2, name: '', threshold: 0.6 },
            'orange': { level: 3, name: '', threshold: 0.8 },
            'red': { level: 4, name: '', threshold: 0.95 }
        };
        
        this.indicators = new Map();
        this.activeWarnings = new Map();
        this.subscriptions = new Map();
        
        this.initializeIndicators();
    }
    
    initializeIndicators() {
        // 
        this.indicators.set('water_level', {
            name: '',
            calculate: this.calculateWaterLevelRisk.bind(this),
            thresholds: {
                warning: 185.0,
                alert: 188.0,
                danger: 190.0,
                emergency: 192.0
            }
        });
        
        // 
        this.indicators.set('rainfall', {
            name: '',
            calculate: this.calculateRainfallRisk.bind(this),
            thresholds: {
                warning: 50,    // 2450mm
                alert: 100,     // 24100mm
                danger: 200,    // 24200mm
                emergency: 300  // 24300mm
            }
        });
        
        // 
        this.indicators.set('dam_safety', {
            name: '',
            calculate: this.calculateDamSafetyRisk.bind(this),
            parameters: ['seepage', 'displacement', 'stress', 'vibration']
        });
    }
    
    async evaluateRiskLevel(stationData, forecastData) {
        const riskScores = new Map();
        
        // 
        for (const [indicatorName, indicator] of this.indicators) {
            const score = await indicator.calculate(stationData, forecastData);
            riskScores.set(indicatorName, score);
        }
        
        // 
        const overallRisk = this.calculateOverallRisk(riskScores);
        
        // 
        const warningLevel = this.determineWarningLevel(overallRisk);
        
        return {
            overallRisk,
            warningLevel,
            indicatorScores: Object.fromEntries(riskScores),
            timestamp: new Date().toISOString()
        };
    }
    
    calculateWaterLevelRisk(stationData, forecastData) {
        const currentLevel = stationData.water_level;
        const thresholds = this.indicators.get('water_level').thresholds;
        
        // 
        let currentRisk = 0;
        if (currentLevel >= thresholds.emergency) currentRisk = 1.0;
        else if (currentLevel >= thresholds.danger) currentRisk = 0.8;
        else if (currentLevel >= thresholds.alert) currentRisk = 0.6;
        else if (currentLevel >= thresholds.warning) currentRisk = 0.3;
        
        // 
        let forecastRisk = 0;
        if (forecastData && forecastData.water_level_forecast) {
            const maxForecastLevel = Math.max(...forecastData.water_level_forecast);
            if (maxForecastLevel >= thresholds.emergency) forecastRisk = 1.0;
            else if (maxForecastLevel >= thresholds.danger) forecastRisk = 0.8;
            else if (maxForecastLevel >= thresholds.alert) forecastRisk = 0.6;
            else if (maxForecastLevel >= thresholds.warning) forecastRisk = 0.3;
        }
        
        // 
        const trendRisk = this.calculateWaterLevelTrend(stationData.recent_data);
        
        return Math.max(currentRisk, forecastRisk, trendRisk);
    }
    
    calculateRainfallRisk(stationData, forecastData) {
        // 24
        const rainfall24h = this.calculateAccumulatedRainfall(stationData.recent_data, 24);
        const thresholds = this.indicators.get('rainfall').thresholds;
        
        let risk = 0;
        if (rainfall24h >= thresholds.emergency) risk = 1.0;
        else if (rainfall24h >= thresholds.danger) risk = 0.8;
        else if (rainfall24h >= thresholds.alert) risk = 0.6;
        else if (rainfall24h >= thresholds.warning) risk = 0.3;
        
        // 
        if (forecastData && forecastData.rainfall_forecast) {
            const forecastTotal = forecastData.rainfall_forecast.reduce((sum, val) => sum + val, 0);
            const combinedRainfall = rainfall24h + forecastTotal;
            
            let forecastRisk = 0;
            if (combinedRainfall >= thresholds.emergency) forecastRisk = 1.0;
            else if (combinedRainfall >= thresholds.danger) forecastRisk = 0.8;
            else if (combinedRainfall >= thresholds.alert) forecastRisk = 0.6;
            else if (combinedRainfall >= thresholds.warning) forecastRisk = 0.3;
            
            risk = Math.max(risk, forecastRisk);
        }
        
        return risk;
    }
    
    async triggerWarning(warningData) {
        const warningId = this.generateWarningId();
        const warning = {
            id: warningId,
            level: warningData.warningLevel,
            type: warningData.type || '',
            stations: warningData.stations,
            risk_score: warningData.overallRisk,
            details: warningData.indicatorScores,
            issued_at: new Date().toISOString(),
            status: 'active'
        };
        
        this.activeWarnings.set(warningId, warning);
        
        // 
        await this.sendWarningNotifications(warning);
        
        // 
        if (warning.level >= 3) {
            await this.activateEmergencyPlan(warning);
        }
        
        return warning;
    }
    
    async sendWarningNotifications(warning) {
        const notifications = [];
        
        // 
        notifications.push(this.sendSMSAlert(warning));
        
        // 
        notifications.push(this.sendEmailAlert(warning));
        
        // 
        notifications.push(this.sendSystemNotification(warning));
        
        // 
        notifications.push(this.sendExternalAPINotification(warning));
        
        await Promise.all(notifications);
    }
}
\end{lstlisting}

\hypertarget{section-106}{%
\subsubsection{}\label{section-106}}

\begin{lstlisting}[style=python]
import numpy as np
import pandas as pd
from sklearn.ensemble import IsolationForest
from sklearn.preprocessing import StandardScaler
import joblib

class IntelligentEventDetector:
    """"""
    
    def __init__(self, config):
        self.config = config
        self.models = {}
        self.scalers = {}
        self.detection_rules = self.load_detection_rules()
        
        # 
        self.load_trained_models()
    
    def load_trained_models(self):
        """"""
        try:
            # 
            self.models['water_level'] = joblib.load('models/water_level_anomaly_model.pkl')
            self.scalers['water_level'] = joblib.load('models/water_level_scaler.pkl')
            
            # 
            self.models['rainfall'] = joblib.load('models/rainfall_anomaly_model.pkl')
            self.scalers['rainfall'] = joblib.load('models/rainfall_scaler.pkl')
            
            # 
            self.models['comprehensive'] = joblib.load('models/comprehensive_anomaly_model.pkl')
            self.scalers['comprehensive'] = joblib.load('models/comprehensive_scaler.pkl')
            
        except FileNotFoundError:
            print("")
            self.train_online_models()
    
    def detect_anomalies(self, data, detection_type='comprehensive'):
        """"""
        
        # 
        processed_data = self.preprocess_data(data, detection_type)
        
        if processed_data.empty:
            return []
        
        # 
        anomalies = []
        
        # 
        statistical_anomalies = self.statistical_anomaly_detection(processed_data)
        anomalies.extend(statistical_anomalies)
        
        # 
        if detection_type in self.models:
            ml_anomalies = self.ml_anomaly_detection(processed_data, detection_type)
            anomalies.extend(ml_anomalies)
        
        # 
        rule_anomalies = self.rule_based_detection(processed_data)
        anomalies.extend(rule_anomalies)
        
        # 
        unique_anomalies = self.deduplicate_anomalies(anomalies)
        sorted_anomalies = sorted(unique_anomalies, key=lambda x: x['severity'], reverse=True)
        
        return sorted_anomalies
    
    def statistical_anomaly_detection(self, data):
        """"""
        anomalies = []
        
        for column in data.select_dtypes(include=[np.number]).columns:
            if column in ['timestamp']:
                continue
                
            values = data[column].dropna()
            if len(values) < 10:  # 
                continue
            
            # Z-score
            z_scores = np.abs((values - values.mean()) / values.std())
            z_anomalies = values[z_scores > 3]
            
            for idx, value in z_anomalies.items():
                anomalies.append({
                    'timestamp': data.loc[idx, 'timestamp'],
                    'parameter': column,
                    'value': value,
                    'anomaly_score': z_scores[idx],
                    'method': 'z_score',
                    'severity': min(z_scores[idx] / 3, 1.0),
                    'description': f'{column}{value}Z-score: {z_scores[idx]:.2f}'
                })
            
            # IQR
            q1 = values.quantile(0.25)
            q3 = values.quantile(0.75)
            iqr = q3 - q1
            lower_bound = q1 - 1.5 \highlight{ iqr
            upper_bound = q3 + 1.5 } iqr
            
            iqr_anomalies = values[(values < lower_bound) | (values > upper_bound)]
            
            for idx, value in iqr_anomalies.items():
                severity = max(
                    (lower_bound - value) / iqr if value < lower_bound else 0,
                    (value - upper_bound) / iqr if value > upper_bound else 0
                ) / 1.5
                
                anomalies.append({
                    'timestamp': data.loc[idx, 'timestamp'],
                    'parameter': column,
                    'value': value,
                    'anomaly_score': severity,
                    'method': 'iqr',
                    'severity': min(severity, 1.0),
                    'description': f'{column}{value}{lower_bound:.2f}-{upper_bound:.2f}'
                })
        
        return anomalies
    
    def ml_anomaly_detection(self, data, detection_type):
        """"""
        anomalies = []
        
        try:
            model = self.models[detection_type]
            scaler = self.scalers[detection_type]
            
            # 
            features = self.extract_features(data, detection_type)
            if features.empty:
                return anomalies
            
            # 
            scaled_features = scaler.transform(features)
            
            # 
            anomaly_scores = model.decision_function(scaled_features)
            predictions = model.predict(scaled_features)
            
            # 
            for i, (score, prediction) in enumerate(zip(anomaly_scores, predictions)):
                if prediction == -1:  # 
                    anomalies.append({
                        'timestamp': data.iloc[i]['timestamp'],
                        'parameter': detection_type,
                        'anomaly_score': abs(score),
                        'method': 'isolation_forest',
                        'severity': min(abs(score), 1.0),
                        'description': f'{detection_type}{score:.3f}'
                    })
        
        except Exception as e:
            print(f"{e}")
        
        return anomalies
    
    def rule_based_detection(self, data):
        """"""
        anomalies = []
        
        for rule in self.detection_rules:
            try:
                # 
                if self.evaluate_rule_condition(data, rule['condition']):
                    anomalies.append({
                        'timestamp': data.iloc[-1]['timestamp'] if not data.empty else pd.Timestamp.now(),
                        'parameter': rule['parameter'],
                        'rule_id': rule['id'],
                        'method': 'rule_based',
                        'severity': rule['severity'],
                        'description': rule['description'],
                        'action': rule.get('action', 'alert')
                    })
            except Exception as e:
                print(f" {rule['id']} {e}")
        
        return anomalies
    
    def evaluate_rule_condition(self, data, condition):
        """"""
        try:
            # 
            # 
            namespace = {
                'data': data,
                'np': np,
                'pd': pd,
                'latest': data.iloc[-1] if not data.empty else None
            }
            
            return eval(condition, {"__builtins__": {}}, namespace)
        except:
            return False
    
    def train_online_models(self):
        """"""
        print("...")
        
        # 
        # 
        
        for detection_type in ['water_level', 'rainfall', 'comprehensive']:
            # 
            training_data = self.generate_training_data(detection_type)
            
            # 
            model = IsolationForest(
                contamination=0.1,
                random_state=42,
                n_estimators=100
            )
            
            scaler = StandardScaler()
            scaled_data = scaler.fit_transform(training_data)
            
            model.fit(scaled_data)
            
            # 
            self.models[detection_type] = model
            self.scalers[detection_type] = scaler
            
            print(f"{detection_type} ")
\end{lstlisting}

\hypertarget{section-107}{%
\subsection{8.4.2}\label{section-107}}

\hypertarget{section-108}{%
\subsubsection{}\label{section-108}}

\begin{lstlisting}[style=javascript]
class EmergencyPlanManager {
    constructor(config) {
        this.config = config;
        this.plans = new Map();
        this.activePlans = new Map();
        this.planHistory = [];
        
        this.loadEmergencyPlans();
    }
    
    loadEmergencyPlans() {
        // 
        this.plans.set('flood', {
            id: 'flood_response',
            name: '',
            trigger_conditions: {
                water_level: { threshold: 188.0, operator: '>=' },
                rainfall_24h: { threshold: 100, operator: '>=' },
                risk_level: { threshold: 0.7, operator: '>=' }
            },
            phases: [
                {
                    phase: 'preparation',
                    name: '',
                    duration: 30, // 
                    actions: [
                        'activate_emergency_center',
                        'notify_key_personnel',
                        'check_communication_systems',
                        'prepare_emergency_supplies'
                    ]
                },
                {
                    phase: 'response',
                    name: '',
                    duration: 120,
                    actions: [
                        'implement_flood_control_measures',
                        'coordinate_evacuation',
                        'monitor_water_levels',
                        'manage_reservoir_operations'
                    ]
                },
                {
                    phase: 'recovery',
                    name: '',
                    duration: 480,
                    actions: [
                        'assess_damage',
                        'restore_normal_operations',
                        'conduct_post_event_analysis',
                        'update_emergency_plans'
                    ]
                }
            ],
            resources: {
                personnel: ['emergency_manager', 'technical_experts', 'operators'],
                equipment: ['pumps', 'generators', 'communication_devices'],
                materials: ['sandbags', 'barriers', 'emergency_supplies']
            }
        });
        
        // 
        this.plans.set('dam_safety', {
            id: 'dam_safety_response',
            name: '',
            trigger_conditions: {
                seepage_rate: { threshold: 0.5, operator: '>=' },
                displacement: { threshold: 10, operator: '>=' },
                structural_alert: { threshold: 'red', operator: '==' }
            },
            phases: [
                {
                    phase: 'immediate',
                    name: '',
                    duration: 10,
                    actions: [
                        'activate_dam_safety_protocol',
                        'notify_dam_safety_team',
                        'implement_emergency_monitoring',
                        'prepare_evacuation_notice'
                    ]
                },
                {
                    phase: 'assessment',
                    name: '',
                    duration: 60,
                    actions: [
                        'conduct_structural_assessment',
                        'analyze_monitoring_data',
                        'determine_risk_level',
                        'decide_mitigation_measures'
                    ]
                },
                {
                    phase: 'mitigation',
                    name: '',
                    duration: 240,
                    actions: [
                        'implement_structural_repairs',
                        'adjust_reservoir_operations',
                        'enhance_monitoring_coverage',
                        'coordinate_downstream_protection'
                    ]
                }
            ]
        });
    }
    
    async activatePlan(planId, triggerData, overrides = {}) {
        const plan = this.plans.get(planId);
        if (!plan) {
            throw new Error(\code{: ${planId}});
        }
        
        // 
        const execution = {
            id: this.generateExecutionId(),
            planId: planId,
            plan: plan,
            triggerData: triggerData,
            overrides: overrides,
            status: 'active',
            currentPhase: 0,
            startTime: new Date(),
            phases: plan.phases.map(phase => ({
                ...phase,
                status: 'pending',
                startTime: null,
                endTime: null,
                actions: phase.actions.map(action => ({
                    id: action,
                    status: 'pending',
                    assignee: null,
                    startTime: null,
                    endTime: null,
                    result: null
                }))
            }))
        };
        
        this.activePlans.set(execution.id, execution);
        
        // 
        await this.startPhase(execution, 0);
        
        // 
        this.logPlanActivation(execution);
        
        return execution;
    }
    
    async startPhase(execution, phaseIndex) {
        if (phaseIndex >= execution.phases.length) {
            await this.completePlan(execution);
            return;
        }
        
        const phase = execution.phases[phaseIndex];
        phase.status = 'active';
        phase.startTime = new Date();
        
        execution.currentPhase = phaseIndex;
        
        console.log(\code{: ${phase.name}});
        
        // 
        const actionPromises = phase.actions.map(action => 
            this.executeAction(execution, phaseIndex, action)
        );
        
        // 
        const timeout = phase.duration \highlight{ 60 } 1000; // 
        
        try {
            await Promise.race([
                Promise.all(actionPromises),
                new Promise((_, reject) => 
                    setTimeout(() => reject(new Error('')), timeout)
                )
            ]);
            
            phase.status = 'completed';
            phase.endTime = new Date();
            
            // 
            await this.startPhase(execution, phaseIndex + 1);
            
        } catch (error) {
            phase.status = 'failed';
            phase.endTime = new Date();
            phase.error = error.message;
            
            // 
            await this.handlePhaseFailure(execution, phaseIndex, error);
        }
    }
    
    async executeAction(execution, phaseIndex, action) {
        action.status = 'running';
        action.startTime = new Date();
        
        try {
            // 
            const result = await this.performAction(action.id, execution);
            
            action.status = 'completed';
            action.endTime = new Date();
            action.result = result;
            
            console.log(\code{: ${action.id}});
            
        } catch (error) {
            action.status = 'failed';
            action.endTime = new Date();
            action.error = error.message;
            
            console.error(\code{: ${action.id} - ${error.message}});
            throw error;
        }
    }
    
    async performAction(actionId, execution) {
        const actionHandlers = {
            // 
            notify_key_personnel: async () => {
                const notifications = await this.sendEmergencyNotifications(execution);
                return { notifications_sent: notifications.length };
            },
            
            // 
            activate_emergency_center: async () => {
                await this.activateEmergencyCenter(execution);
                return { center_activated: true };
            },
            
            // 
            implement_flood_control_measures: async () => {
                const measures = await this.implementFloodControlMeasures(execution);
                return { measures_implemented: measures };
            },
            
            // 
            monitor_water_levels: async () => {
                await this.enhanceWaterLevelMonitoring(execution);
                return { enhanced_monitoring: true };
            },
            
            // 
            assess_damage: async () => {
                const assessment = await this.conductDamageAssessment(execution);
                return assessment;
            }
        };
        
        const handler = actionHandlers[actionId];
        if (handler) {
            return await handler();
        } else {
            throw new Error(\code{: ${actionId}});
        }
    }
    
    async sendEmergencyNotifications(execution) {
        const notifications = [];
        const personnel = execution.plan.resources.personnel;
        
        for (const role of personnel) {
            const contacts = await this.getPersonnelContacts(role);
            
            for (const contact of contacts) {
                try {
                    await this.sendNotification(contact, {
                        type: 'emergency_activation',
                        plan: execution.plan.name,
                        trigger: execution.triggerData,
                        urgency: 'high'
                    });
                    
                    notifications.push({
                        recipient: contact.name,
                        method: contact.preferred_method,
                        status: 'sent'
                    });
                } catch (error) {
                    notifications.push({
                        recipient: contact.name,
                        method: contact.preferred_method,
                        status: 'failed',
                        error: error.message
                    });
                }
            }
        }
        
        return notifications;
    }
    
    getExecutionStatus(executionId) {
        const execution = this.activePlans.get(executionId);
        if (!execution) {
            return null;
        }
        
        const currentPhase = execution.phases[execution.currentPhase];
        const completedActions = currentPhase ? 
            currentPhase.actions.filter(a => a.status === 'completed').length : 0;
        const totalActions = currentPhase ? currentPhase.actions.length : 0;
        
        return {
            executionId: execution.id,
            planName: execution.plan.name,
            status: execution.status,
            currentPhase: currentPhase ? currentPhase.name : null,
            progress: totalActions > 0 ? (completedActions / totalActions \highlight{ 100) : 0,
            startTime: execution.startTime,
            estimatedCompletion: this.calculateEstimatedCompletion(execution)
        };
    }
}
\end{lstlisting}

\hypertarget{section-109}{%
\subsection{8.4.3}\label{section-109}}

\hypertarget{section-110}{%
\subsubsection{}\label{section-110}}

\begin{lstlisting}[style=python]
class SituationAnalysisEngine:
    """"""
    
    def __init__(self, config):
        self.config = config
        self.analysis_modules = {
            'flood_risk': FloodRiskAnalyzer(),
            'dam_safety': DamSafetyAnalyzer(), 
            'water_supply': WaterSupplyAnalyzer(),
            'drought_risk': DroughtRiskAnalyzer()
        }
        
    def analyze_current_situation(self, data_snapshot):
        """"""
        
        analysis_results = {}
        
        # 
        for module_name, analyzer in self.analysis_modules.items():
            try:
                result = analyzer.analyze(data_snapshot)
                analysis_results[module_name] = result
            except Exception as e:
                print(f" {module_name} : {e}")
                analysis_results[module_name] = {
                    'status': 'error',
                    'error': str(e)
                }
        
        # 
        overall_assessment = self.generate_overall_assessment(analysis_results)
        
        return {
            'timestamp': data_snapshot['timestamp'],
            'module_results': analysis_results,
            'overall_assessment': overall_assessment,
            'recommendations': self.generate_recommendations(analysis_results)
        }
    
    def generate_overall_assessment(self, analysis_results):
        """"""
        
        # 
        risk_levels = []
        for module_name, result in analysis_results.items():
            if result.get('status') == 'success' and 'risk_level' in result:
                risk_levels.append(result['risk_level'])
        
        if not risk_levels:
            return {
                'overall_risk': 'unknown',
                'confidence': 0,
                'summary': ''
            }
        
        # 
        max_risk = max(risk_levels)
        avg_risk = sum(risk_levels) / len(risk_levels)
        
        # 
        if max_risk >= 0.8:
            overall_risk = 'critical'
        elif max_risk >= 0.6:
            overall_risk = 'high'
        elif avg_risk >= 0.4:
            overall_risk = 'medium'
        else:
            overall_risk = 'low'
        
        # 
        confidence = min(len(risk_levels) / len(self.analysis_modules), 1.0)
        
        return {
            'overall_risk': overall_risk,
            'max_risk_level': max_risk,
            'average_risk_level': avg_risk,
            'confidence': confidence,
            'summary': self.generate_risk_summary(overall_risk, analysis_results)
        }
    
    def generate_recommendations(self, analysis_results):
        """"""
        
        recommendations = []
        
        for module_name, result in analysis_results.items():
            if result.get('status') == 'success' and 'recommendations' in result:
                module_recommendations = result['recommendations']
                
                for rec in module_recommendations:
                    recommendations.append({
                        'source_module': module_name,
                        'priority': rec.get('priority', 'medium'),
                        'action': rec['action'],
                        'rationale': rec.get('rationale', ''),
                        'estimated_impact': rec.get('impact', 'unknown'),
                        'implementation_time': rec.get('time_required', 'unknown')
                    })
        
        # 
        priority_order = {'critical': 4, 'high': 3, 'medium': 2, 'low': 1}
        recommendations.sort(
            key=lambda x: priority_order.get(x['priority'], 0),
            reverse=True
        )
        
        return recommendations

class FloodRiskAnalyzer:
    """"""
    
    def analyze(self, data_snapshot):
        """"""
        
        try:
            # 
            water_levels = data_snapshot.get('water_levels', {})
            rainfall_data = data_snapshot.get('rainfall', {})
            weather_forecast = data_snapshot.get('weather_forecast', {})
            
            # 
            water_level_risk = self.assess_water_level_risk(water_levels)
            rainfall_risk = self.assess_rainfall_risk(rainfall_data)
            forecast_risk = self.assess_forecast_risk(weather_forecast)
            
            # 
            combined_risk = max(water_level_risk, rainfall_risk, forecast_risk)
            
            # 
            recommendations = self.generate_flood_recommendations(
                combined_risk, water_level_risk, rainfall_risk, forecast_risk
            )
            
            return {
                'status': 'success',
                'risk_level': combined_risk,
                'factors': {
                    'water_level_risk': water_level_risk,
                    'rainfall_risk': rainfall_risk,
                    'forecast_risk': forecast_risk
                },
                'recommendations': recommendations,
                'details': {
                    'critical_stations': self.identify_critical_stations(water_levels),
                    'peak_forecast': self.predict_flood_peak(data_snapshot)
                }
            }
            
        except Exception as e:
            return {
                'status': 'error',
                'error': str(e)
            }
    
    def assess_water_level_risk(self, water_levels):
        """"""
        if not water_levels:
            return 0
        
        max_risk = 0
        for station_id, level_data in water_levels.items():
            current_level = level_data.get('current_level', 0)
            warning_level = level_data.get('warning_level', float('inf'))
            alert_level = level_data.get('alert_level', float('inf'))
            
            if current_level >= alert_level:
                risk = 0.9
            elif current_level >= warning_level:
                risk = 0.6
            else:
                # 
                risk = max(0, (current_level - warning_level } 0.8) / (warning_level * 0.2))
            
            max_risk = max(max_risk, risk)
        
        return min(max_risk, 1.0)
    
    def generate_flood_recommendations(self, combined_risk, water_risk, rainfall_risk, forecast_risk):
        """"""
        
        recommendations = []
        
        if combined_risk >= 0.8:
            recommendations.extend([
                {
                    'priority': 'critical',
                    'action': '',
                    'rationale': f' {combined_risk:.2f}',
                    'impact': 'high',
                    'time_required': ''
                },
                {
                    'priority': 'critical', 
                    'action': '',
                    'rationale': '',
                    'impact': 'critical',
                    'time_required': '30'
                }
            ])
        
        if water_risk >= 0.6:
            recommendations.append({
                'priority': 'high',
                'action': '',
                'rationale': f' {water_risk:.2f}',
                'impact': 'high',
                'time_required': '1'
            })
        
        if rainfall_risk >= 0.5:
            recommendations.append({
                'priority': 'medium',
                'action': '',
                'rationale': f' {rainfall_risk:.2f}',
                'impact': 'medium',
                'time_required': ''
            })
        
        return recommendations
\end{lstlisting}

\hypertarget{section-111}{%
\subsection{}\label{section-111}}

\important{}

\begin{enumerate}
\def\labelenumi{\arabic{enumi}.}
\item
  \important{}
\item
  \important{}
\item
  \important{}
\item
  \important{}
\end{enumerate}

\hypertarget{section-112}{%
\subsection{}\label{section-112}}

\hypertarget{section-113}{%
\subsubsection{}\label{section-113}}

\begin{enumerate}
\def\labelenumi{\arabic{enumi}.}
\tightlist
\item
\item
\item
\end{enumerate}

\hypertarget{section-114}{%
\subsubsection{}\label{section-114}}

\begin{enumerate}
\def\labelenumi{\arabic{enumi}.}
\tightlist
\item
\item
\item
\end{enumerate}

\hypertarget{section-115}{%
\subsubsection{}\label{section-115}}

\begin{enumerate}
\def\labelenumi{\arabic{enumi}.}
\tightlist
\item
\item
\end{enumerate}

\hypertarget{section-116}{%
\subsection{}\label{section-116}}

\hypertarget{section-117}{%
\section{}\label{section-117}}

\hypertarget{section-118}{%
\subsection{}\label{section-118}}

\hypertarget{section-119}{%
\subsection{8.5.1}\label{section-119}}

\hypertarget{section-120}{%
\subsubsection{}\label{section-120}}

\begin{lstlisting}[style=javascript]
// 
class ServiceRegistry {
    constructor(config) {
        this.services = new Map();
        this.healthChecks = new Map();
        this.loadBalancers = new Map();
        
        this.config = config;
        this.initializeRegistry();
    }
    
    initializeRegistry() {
        // 
        this.registerService('data-collection', {
            instances: [
                { id: 'dc-001', host: '10.1.1.10', port: 8080, weight: 100 },
                { id: 'dc-002', host: '10.1.1.11', port: 8080, weight: 100 }
            ],
            healthCheck: {
                path: '/health',
                interval: 30000,
                timeout: 5000
            }
        });
        
        this.registerService('data-processing', {
            instances: [
                { id: 'dp-001', host: '10.1.2.10', port: 8081, weight: 100 },
                { id: 'dp-002', host: '10.1.2.11', port: 8081, weight: 100 },
                { id: 'dp-003', host: '10.1.2.12', port: 8081, weight: 100 }
            ],
            healthCheck: {
                path: '/health',
                interval: 15000,
                timeout: 3000
            }
        });
        
        this.registerService('visualization', {
            instances: [
                { id: 'viz-001', host: '10.1.3.10', port: 8082, weight: 100 },
                { id: 'viz-002', host: '10.1.3.11', port: 8082, weight: 100 }
            ],
            healthCheck: {
                path: '/health',
                interval: 30000,
                timeout: 5000
            }
        });
    }
    
    registerService(serviceName, config) {
        this.services.set(serviceName, {
            name: serviceName,
            instances: config.instances,
            healthCheck: config.healthCheck,
            status: 'active',
            registeredAt: new Date()
        });
        
        // 
        this.loadBalancers.set(serviceName, new LoadBalancer(config.instances));
        
        // 
        this.startHealthCheck(serviceName, config.healthCheck);
    }
    
    async startHealthCheck(serviceName, healthConfig) {
        const service = this.services.get(serviceName);
        
        const healthCheck = setInterval(async () => {
            for (const instance of service.instances) {
                try {
                    const response = await fetch(
                        \code{http://${instance.host}:${instance.port}${healthConfig.path}},
                        { timeout: healthConfig.timeout }
                    );
                    
                    instance.healthy = response.ok;
                    instance.lastCheck = new Date();
                    
                    if (!response.ok) {
                        console.warn(\code{ ${instance.id} });
                    }
                } catch (error) {
                    instance.healthy = false;
                    instance.lastCheck = new Date();
                    instance.lastError = error.message;
                    
                    console.error(\code{ ${instance.id} : ${error.message}});
                }
            }
        }, healthConfig.interval);
        
        this.healthChecks.set(serviceName, healthCheck);
    }
    
    getServiceInstance(serviceName) {
        const loadBalancer = this.loadBalancers.get(serviceName);
        return loadBalancer ? loadBalancer.getNextInstance() : null;
    }
    
    getServiceStatus() {
        const status = {};
        
        for (const [serviceName, service] of this.services) {
            const healthyInstances = service.instances.filter(i => i.healthy).length;
            const totalInstances = service.instances.length;
            
            status[serviceName] = {
                total: totalInstances,
                healthy: healthyInstances,
                ratio: totalInstances > 0 ? (healthyInstances / totalInstances) : 0,
                instances: service.instances.map(i => ({
                    id: i.id,
                    host: i.host,
                    port: i.port,
                    healthy: i.healthy,
                    lastCheck: i.lastCheck
                }))
            };
        }
        
        return status;
    }
}

// 
class LoadBalancer {
    constructor(instances) {
        this.instances = instances;
        this.algorithm = 'round_robin';
        this.currentIndex = 0;
    }
    
    getNextInstance() {
        const healthyInstances = this.instances.filter(i => i.healthy);
        
        if (healthyInstances.length === 0) {
            throw new Error('');
        }
        
        switch (this.algorithm) {
            case 'round_robin':
                return this.roundRobin(healthyInstances);
            case 'weighted_round_robin':
                return this.weightedRoundRobin(healthyInstances);
            case 'least_connections':
                return this.leastConnections(healthyInstances);
            default:
                return healthyInstances[0];
        }
    }
    
    roundRobin(instances) {
        const instance = instances[this.currentIndex % instances.length];
        this.currentIndex++;
        return instance;
    }
    
    weightedRoundRobin(instances) {
        // 
        const totalWeight = instances.reduce((sum, i) => sum + i.weight, 0);
        const random = Math.random() * totalWeight;
        
        let currentWeight = 0;
        for (const instance of instances) {
            currentWeight += instance.weight;
            if (random <= currentWeight) {
                return instance;
            }
        }
        
        return instances[0];
    }
}
\end{lstlisting}

\hypertarget{section-121}{%
\subsubsection{}\label{section-121}}

\begin{lstlisting}
#  (istio)
apiVersion: networking.istio.io/v1alpha3
kind: VirtualService
metadata:
  name: water-monitoring-platform
  namespace: smart-water
spec:
  http:
  - match:
    - uri:
        prefix: /api/data-collection
    route:
    - destination:
        host: data-collection-service
        subset: v1
      weight: 90
    - destination:
        host: data-collection-service
        subset: v2
      weight: 10
    fault:
      delay:
        percentage:
          value: 0.1
        fixedDelay: 5s
    retries:
      attempts: 3
      perTryTimeout: 2s
      
  - match:
    - uri:
        prefix: /api/processing
    route:
    - destination:
        host: data-processing-service
    circuitBreaker:
      consecutiveErrors: 5
      interval: 30s
      baseEjectionTime: 30s
      
  - match:
    - uri:
        prefix: /api/visualization
    route:
    - destination:
        host: visualization-service
    timeout: 30s
    mirror:
      host: visualization-service-canary
      
---
apiVersion: networking.istio.io/v1alpha3
kind: DestinationRule
metadata:
  name: data-collection-destination
  namespace: smart-water
spec:
  host: data-collection-service
  trafficPolicy:
    connectionPool:
      tcp:
        maxConnections: 100
      http:
        http1MaxPendingRequests: 50
        maxRequestsPerConnection: 10
    loadBalancer:
      simple: LEAST_CONN
    outlierDetection:
      consecutiveErrors: 3
      interval: 30s
      baseEjectionTime: 30s
  subsets:
  - name: v1
    labels:
      version: v1
  - name: v2
    labels:
      version: v2
\end{lstlisting}

\hypertarget{section-122}{%
\subsection{8.5.2}\label{section-122}}

\hypertarget{section-123}{%
\subsubsection{}\label{section-123}}

\begin{lstlisting}[style=python]
import psutil
import time
from datetime import datetime
import json
import requests
from dataclasses import dataclass
from typing import Dict, List, Optional

@dataclass
class MetricPoint:
    timestamp: float
    value: float
    tags: Dict[str, str] = None

class SystemMonitor:
    """"""
    
    def __init__(self, config):
        self.config = config
        self.metrics_buffer = []
        self.alert_rules = self.load_alert_rules()
        self.collectors = self.initialize_collectors()
        
    def initialize_collectors(self):
        """"""
        return {
            'system': SystemMetricsCollector(),
            'application': ApplicationMetricsCollector(),
            'business': BusinessMetricsCollector(),
            'network': NetworkMetricsCollector()
        }
    
    def collect_all_metrics(self):
        """"""
        all_metrics = {}
        
        for collector_name, collector in self.collectors.items():
            try:
                metrics = collector.collect()
                all_metrics[collector_name] = metrics
            except Exception as e:
                print(f" {collector_name} : {e}")
                
        return all_metrics
    
    def process_metrics(self, metrics):
        """"""
        processed_metrics = []
        
        for category, category_metrics in metrics.items():
            for metric_name, metric_data in category_metrics.items():
                # 
                point = MetricPoint(
                    timestamp=time.time(),
                    value=metric_data['value'],
                    tags={
                        'category': category,
                        'metric': metric_name,
                        'host': self.config.get('hostname', 'unknown'),
                        **metric_data.get('tags', {})
                    }
                )
                processed_metrics.append(point)
        
        return processed_metrics
    
    def check_alerts(self, metrics):
        """"""
        alerts = []
        
        for rule in self.alert_rules:
            try:
                if self.evaluate_alert_rule(rule, metrics):
                    alert = {
                        'rule_id': rule['id'],
                        'severity': rule['severity'],
                        'message': rule['message'],
                        'timestamp': datetime.now().isoformat(),
                        'metrics': self.extract_relevant_metrics(rule, metrics)
                    }
                    alerts.append(alert)
            except Exception as e:
                print(f" {rule['id']} : {e}")
        
        return alerts
    
    def evaluate_alert_rule(self, rule, metrics):
        """"""
        target_metrics = self.extract_relevant_metrics(rule, metrics)
        
        if not target_metrics:
            return False
        
        condition = rule['condition']
        threshold = rule['threshold']
        
        for metric in target_metrics:
            value = metric.value
            
            if condition == 'greater_than' and value > threshold:
                return True
            elif condition == 'less_than' and value < threshold:
                return True
            elif condition == 'equals' and value == threshold:
                return True
            elif condition == 'not_equals' and value != threshold:
                return True
        
        return False

class SystemMetricsCollector:
    """"""
    
    def collect(self):
        """"""
        return {
            'cpu_usage': {
                'value': psutil.cpu_percent(interval=1),
                'unit': 'percent',
                'tags': {'type': 'system'}
            },
            'memory_usage': {
                'value': psutil.virtual_memory().percent,
                'unit': 'percent',
                'tags': {'type': 'system'}
            },
            'disk_usage': {
                'value': psutil.disk_usage('/').percent,
                'unit': 'percent',
                'tags': {'type': 'system', 'mount_point': '/'}
            },
            'load_average': {
                'value': psutil.getloadavg()[0],
                'unit': 'load',
                'tags': {'type': 'system', 'period': '1min'}
            },
            'network_io_bytes_sent': {
                'value': psutil.net_io_counters().bytes_sent,
                'unit': 'bytes',
                'tags': {'type': 'network', 'direction': 'sent'}
            },
            'network_io_bytes_recv': {
                'value': psutil.net_io_counters().bytes_recv,
                'unit': 'bytes',
                'tags': {'type': 'network', 'direction': 'received'}
            }
        }

class ApplicationMetricsCollector:
    """"""
    
    def __init__(self):
        self.service_endpoints = {
            'data-collection': 'http://localhost:8080/metrics',
            'data-processing': 'http://localhost:8081/metrics',
            'visualization': 'http://localhost:8082/metrics'
        }
    
    def collect(self):
        """"""
        metrics = {}
        
        for service_name, endpoint in self.service_endpoints.items():
            try:
                service_metrics = self.collect_service_metrics(service_name, endpoint)
                metrics.update(service_metrics)
            except Exception as e:
                print(f" {service_name} : {e}")
                # 
                metrics[f'{service_name}_available'] = {
                    'value': 0,
                    'unit': 'boolean',
                    'tags': {'service': service_name, 'status': 'unavailable'}
                }
        
        return metrics
    
    def collect_service_metrics(self, service_name, endpoint):
        """"""
        response = requests.get(endpoint, timeout=5)
        response.raise_for_status()
        
        service_data = response.json()
        
        return {
            f'{service_name}_response_time': {
                'value': service_data.get('avg_response_time', 0),
                'unit': 'milliseconds',
                'tags': {'service': service_name}
            },
            f'{service_name}_request_rate': {
                'value': service_data.get('requests_per_second', 0),
                'unit': 'rps',
                'tags': {'service': service_name}
            },
            f'{service_name}_error_rate': {
                'value': service_data.get('error_rate', 0),
                'unit': 'percent',
                'tags': {'service': service_name}
            },
            f'{service_name}_active_connections': {
                'value': service_data.get('active_connections', 0),
                'unit': 'count',
                'tags': {'service': service_name}
            }
        }

class BusinessMetricsCollector:
    """"""
    
    def __init__(self):
        self.database_connection = self.get_database_connection()
    
    def collect(self):
        """"""
        return {
            'active_monitoring_stations': {
                'value': self.count_active_stations(),
                'unit': 'count',
                'tags': {'type': 'business'}
            },
            'data_points_per_minute': {
                'value': self.count_recent_data_points(),
                'unit': 'count/min',
                'tags': {'type': 'business'}
            },
            'active_alerts': {
                'value': self.count_active_alerts(),
                'unit': 'count',
                'tags': {'type': 'business'}
            },
            'system_users_online': {
                'value': self.count_online_users(),
                'unit': 'count',
                'tags': {'type': 'business'}
            }
        }
    
    def count_active_stations(self):
        """"""
        query = """
        SELECT COUNT(DISTINCT station_id) 
        FROM monitoring_data 
        WHERE timestamp > NOW() - INTERVAL 1 HOUR
        """
        result = self.database_connection.execute(query).fetchone()
        return result[0] if result else 0
    
    def count_recent_data_points(self):
        """"""
        query = """
        SELECT COUNT(\highlight{) 
        FROM monitoring_data 
        WHERE timestamp > NOW() - INTERVAL 1 MINUTE
        """
        result = self.database_connection.execute(query).fetchone()
        return result[0] if result else 0
\end{lstlisting}

\hypertarget{section-124}{%
\subsubsection{}\label{section-124}}

\begin{lstlisting}[style=javascript]
class LogManagementSystem {
    constructor(config) {
        this.config = config;
        this.logLevel = config.logLevel || 'INFO';
        this.logDestinations = config.destinations || [];
        this.logBuffer = [];
        this.maxBufferSize = config.maxBufferSize || 1000;
        
        this.initializeDestinations();
        this.setupPeriodicFlush();
    }
    
    initializeDestinations() {
        this.destinations = this.logDestinations.map(dest => {
            switch (dest.type) {
                case 'file':
                    return new FileLogDestination(dest.config);
                case 'elasticsearch':
                    return new ElasticsearchLogDestination(dest.config);
                case 'kafka':
                    return new KafkaLogDestination(dest.config);
                default:
                    return new ConsoleLogDestination();
            }
        });
    }
    
    log(level, message, metadata = {}) {
        if (!this.shouldLog(level)) {
            return;
        }
        
        const logEntry = {
            timestamp: new Date().toISOString(),
            level: level,
            message: message,
            metadata: {
                ...metadata,
                hostname: this.config.hostname,
                service: this.config.serviceName,
                version: this.config.version,
                pid: process.pid
            },
            traceId: this.getTraceId()
        };
        
        this.logBuffer.push(logEntry);
        
        // 
        if (level === 'ERROR' || level === 'FATAL') {
            this.flush();
        }
        
        // 
        if (this.logBuffer.length >= this.maxBufferSize) {
            this.flush();
        }
    }
    
    shouldLog(level) {
        const levels = ['DEBUG', 'INFO', 'WARN', 'ERROR', 'FATAL'];
        const currentLevelIndex = levels.indexOf(this.logLevel);
        const messageLevelIndex = levels.indexOf(level);
        
        return messageLevelIndex >= currentLevelIndex;
    }
    
    flush() {
        if (this.logBuffer.length === 0) {
            return;
        }
        
        const logs = [...this.logBuffer];
        this.logBuffer = [];
        
        // 
        this.destinations.forEach(destination => {
            destination.send(logs).catch(error => {
                console.error(\code{: ${error.message}});
            });
        });
    }
    
    setupPeriodicFlush() {
        setInterval(() => {
            this.flush();
        }, 5000); // 5
    }
    
    // 
    debug(message, metadata) { this.log('DEBUG', message, metadata); }
    info(message, metadata) { this.log('INFO', message, metadata); }
    warn(message, metadata) { this.log('WARN', message, metadata); }
    error(message, metadata) { this.log('ERROR', message, metadata); }
    fatal(message, metadata) { this.log('FATAL', message, metadata); }
}

class ElasticsearchLogDestination {
    constructor(config) {
        this.config = config;
        this.client = new ElasticsearchClient(config);
        this.indexPattern = config.indexPattern || 'logs-{yyyy.MM.dd}';
    }
    
    async send(logs) {
        const bulkBody = [];
        
        logs.forEach(log => {
            const index = this.generateIndexName(log.timestamp);
            
            bulkBody.push({
                index: { _index: index }
            });
            bulkBody.push(log);
        });
        
        try {
            await this.client.bulk({ body: bulkBody });
        } catch (error) {
            throw new Error(\code{Elasticsearch: ${error.message}});
        }
    }
    
    generateIndexName(timestamp) {
        const date = new Date(timestamp);
        return this.indexPattern
            .replace('{yyyy}', date.getFullYear())
            .replace('{MM}', String(date.getMonth() + 1).padStart(2, '0'))
            .replace('{dd}', String(date.getDate()).padStart(2, '0'));
    }
}
\end{lstlisting}

\hypertarget{section-125}{%
\subsection{8.5.3}\label{section-125}}

\hypertarget{section-126}{%
\subsubsection{}\label{section-126}}

\begin{lstlisting}
# Docker Compose
version: '3.8'

services:
  # 
  data-collection:
    image: smart-water/data-collection:latest
    container_name: data-collection
    ports:
      - "8080:8080"
    environment:
      - SPRING_PROFILES_ACTIVE=production
      - DATABASE_URL=jdbc:postgresql://postgres:5432/smart_water
      - REDIS_URL=redis://redis:6379
    depends_on:
      - postgres
      - redis
    volumes:
      - ./config/data-collection.yml:/app/config/application.yml
      - ./logs:/app/logs
    networks:
      - smart-water-network
    restart: unless-stopped
    healthcheck:
      test: ["CMD", "curl", "-f", "http://localhost:8080/health"]
      interval: 30s
      timeout: 10s
      retries: 3
    deploy:
      resources:
        limits:
          cpus: '1.0'
          memory: 1G
        reservations:
          cpus: '0.5'
          memory: 512M
  
  # 
  data-processing:
    image: smart-water/data-processing:latest
    container_name: data-processing
    ports:
      - "8081:8081"
    environment:
      - NODE_ENV=production
      - DATABASE_URL=postgresql://postgres:5432/smart_water
      - KAFKA_BROKERS=kafka:9092
    depends_on:
      - postgres
      - kafka
    volumes:
      - ./config/data-processing.json:/app/config/production.json
      - ./logs:/app/logs
    networks:
      - smart-water-network
    restart: unless-stopped
    scale: 3
    
  # 
  visualization:
    image: smart-water/visualization:latest
    container_name: visualization
    ports:
      - "8082:8082"
    environment:
      - VUE_APP_API_BASE_URL=http://api-gateway:8000
      - VUE_APP_CESIUM_TOKEN=${CESIUM_ACCESS_TOKEN}
    depends_on:
      - api-gateway
    volumes:
      - ./dist:/app/dist
      - ./nginx.conf:/etc/nginx/nginx.conf
    networks:
      - smart-water-network
    restart: unless-stopped
    
  # API
  api-gateway:
    image: nginx:alpine
    container_name: api-gateway
    ports:
      - "8000:80"
      - "8443:443"
    volumes:
      - ./nginx/gateway.conf:/etc/nginx/conf.d/default.conf
      - ./ssl:/etc/nginx/ssl
    depends_on:
      - data-collection
      - data-processing
      - visualization
    networks:
      - smart-water-network
    restart: unless-stopped
    
  # 
  postgres:
    image: postgis/postgis:13-3.1
    container_name: postgres
    environment:
      - POSTGRES_DB=smart_water
      - POSTGRES_USER=smart_water_user
      - POSTGRES_PASSWORD=${POSTGRES_PASSWORD}
    volumes:
      - postgres_data:/var/lib/postgresql/data
      - ./sql/init.sql:/docker-entrypoint-initdb.d/init.sql
    networks:
      - smart-water-network
    restart: unless-stopped
    
  # Redis
  redis:
    image: redis:6-alpine
    container_name: redis
    command: redis-server --appendonly yes --requirepass ${REDIS_PASSWORD}
    volumes:
      - redis_data:/data
    networks:
      - smart-water-network
    restart: unless-stopped
    
  # 
  prometheus:
    image: prom/prometheus:latest
    container_name: prometheus
    ports:
      - "9090:9090"
    volumes:
      - ./monitoring/prometheus.yml:/etc/prometheus/prometheus.yml
      - prometheus_data:/prometheus
    networks:
      - smart-water-network
    restart: unless-stopped
    
  grafana:
    image: grafana/grafana:latest
    container_name: grafana
    ports:
      - "3000:3000"
    environment:
      - GF_SECURITY_ADMIN_PASSWORD=${GRAFANA_PASSWORD}
    volumes:
      - grafana_data:/var/lib/grafana
      - ./monitoring/grafana/dashboards:/etc/grafana/provisioning/dashboards
    networks:
      - smart-water-network
    restart: unless-stopped

volumes:
  postgres_data:
  redis_data:
  prometheus_data:
  grafana_data:

networks:
  smart-water-network:
    driver: bridge
\end{lstlisting}

\hypertarget{cicd}{%
\subsubsection{CI/CD}\label{cicd}}

\begin{lstlisting}
# GitLab CI/CD
stages:
  - validate
  - test
  - build
  - security-scan
  - deploy-staging
  - integration-test
  - deploy-production
  - post-deploy

variables:
  DOCKER_REGISTRY: registry.smart-water.com
  KUBECONFIG: /tmp/.kube/config

# 
code-quality:
  stage: validate
  image: sonarqube/sonar-scanner-cli:latest
  script:
    - sonar-scanner -Dsonar.projectKey=smart-water-platform
  only:
    - merge_requests
    - develop
    - master

# 
unit-tests:
  stage: test
  image: node:16
  services:
    - postgres:13
    - redis:6
  script:
    - npm ci
    - npm run test:unit
    - npm run test:coverage
  coverage: '/Lines\s}:\s\highlight{(\d+\.?\d})%/'
  artifacts:
    reports:
      coverage_report:
        coverage_format: cobertura
        path: coverage/cobertura-coverage.xml
    expire_in: 1 week

# Docker
build-images:
  stage: build
  image: docker:20.10.16
  services:
    - docker:20.10.16-dind
  before_script:
    - echo $CI_REGISTRY_PASSWORD | docker login -u $CI_REGISTRY_USER --password-stdin $CI_REGISTRY
  script:
    - docker build -t $DOCKER_REGISTRY/data-collection:$CI_COMMIT_SHA ./services/data-collection
    - docker build -t $DOCKER_REGISTRY/data-processing:$CI_COMMIT_SHA ./services/data-processing
    - docker build -t $DOCKER_REGISTRY/visualization:$CI_COMMIT_SHA ./services/visualization
    - docker push $DOCKER_REGISTRY/data-collection:$CI_COMMIT_SHA
    - docker push $DOCKER_REGISTRY/data-processing:$CI_COMMIT_SHA
    - docker push $DOCKER_REGISTRY/visualization:$CI_COMMIT_SHA
  only:
    - develop
    - master

# 
security-scan:
  stage: security-scan
  image: owasp/zap2docker-stable
  script:
    - zap-baseline.py -t http://staging.smart-water.com
  artifacts:
    reports:
      junit: zap-report.xml
  only:
    - develop
    - master

# 
deploy-staging:
  stage: deploy-staging
  image: bitnami/kubectl:latest
  environment:
    name: staging
    url: https://staging.smart-water.com
  script:
    - kubectl config use-context staging
    - helm upgrade --install smart-water-staging ./helm-chart 
        --set image.tag=$CI_COMMIT_SHA
        --set environment=staging
        --namespace smart-water-staging
    - kubectl rollout status deployment/smart-water-staging -n smart-water-staging
  only:
    - develop

# 
deploy-production:
  stage: deploy-production
  image: bitnami/kubectl:latest
  environment:
    name: production
    url: https://smart-water.com
  script:
    - kubectl config use-context production
    - helm upgrade --install smart-water-production ./helm-chart 
        --set image.tag=$CI_COMMIT_SHA
        --set environment=production
        --namespace smart-water-production
    - kubectl rollout status deployment/smart-water-production -n smart-water-production
  when: manual
  only:
    - master

# 
post-deploy-check:
  stage: post-deploy
  image: alpine/curl
  script:
    - curl -f https://smart-water.com/health || exit 1
    - curl -f https://smart-water.com/api/health || exit 1
  only:
    - master
\end{lstlisting}

\hypertarget{section-127}{%
\subsection{8.5.4}\label{section-127}}

\hypertarget{section-128}{%
\subsubsection{}\label{section-128}}

\begin{lstlisting}[style=python]
class AutoRecoverySystem:
    """"""
    
    def __init__(self, config):
        self.config = config
        self.recovery_strategies = self.load_recovery_strategies()
        self.circuit_breakers = {}
        self.retry_policies = {}
        
    def load_recovery_strategies(self):
        """"""
        return {
            'service_unavailable': {
                'detection': {
                    'method': 'health_check_failure',
                    'threshold': 3,
                    'window': 60  # 
                },
                'recovery': [
                    'restart_service',
                    'scale_out_replicas',
                    'fallback_to_backup'
                ]
            },
            'high_cpu_usage': {
                'detection': {
                    'method': 'metric_threshold',
                    'metric': 'cpu_usage',
                    'threshold': 80,
                    'duration': 300
                },
                'recovery': [
                    'scale_out_replicas',
                    'optimize_resource_allocation',
                    'enable_performance_mode'
                ]
            },
            'memory_leak': {
                'detection': {
                    'method': 'memory_growth_trend',
                    'growth_rate': 10,  # MB/minute
                    'duration': 600
                },
                'recovery': [
                    'restart_service_gracefully',
                    'collect_heap_dump',
                    'notify_development_team'
                ]
            },
            'database_connection_pool_exhausted': {
                'detection': {
                    'method': 'connection_pool_metric',
                    'threshold': 95,  # percent
                    'duration': 60
                },
                'recovery': [
                    'increase_connection_pool_size',
                    'kill_long_running_queries',
                    'restart_database_connections'
                ]
            }
        }
    
    async def monitor_and_recover(self):
        """"""
        while True:
            try:
                # 
                metrics = await self.collect_system_metrics()
                
                # 
                detected_issues = self.detect_issues(metrics)
                
                # 
                for issue in detected_issues:
                    await self.execute_recovery(issue)
                
                await asyncio.sleep(30)  # 30
                
            except Exception as e:
                print(f": {e}")
                await asyncio.sleep(60)
    
    def detect_issues(self, metrics):
        """"""
        issues = []
        
        for strategy_name, strategy in self.recovery_strategies.items():
            if self.evaluate_detection_rule(strategy['detection'], metrics):
                issues.append({
                    'type': strategy_name,
                    'strategy': strategy,
                    'detected_at': datetime.now(),
                    'metrics': metrics
                })
        
        return issues
    
    async def execute_recovery(self, issue):
        """"""
        strategy = issue['strategy']
        
        print(f": {issue['type']}, ")
        
        for recovery_action in strategy['recovery']:
            try:
                success = await self.execute_recovery_action(recovery_action, issue)
                
                if success:
                    print(f": {recovery_action}")
                    # 
                    if await self.verify_recovery(issue):
                        print(f": {issue['type']}")
                        break
                else:
                    print(f": {recovery_action}")
                    
            except Exception as e:
                print(f" {recovery_action}: {e}")
    
    async def execute_recovery_action(self, action, issue):
        """"""
        
        action_handlers = {
            'restart_service': self.restart_service,
            'scale_out_replicas': self.scale_out_replicas,
            'fallback_to_backup': self.fallback_to_backup,
            'optimize_resource_allocation': self.optimize_resource_allocation,
            'increase_connection_pool_size': self.increase_connection_pool_size,
            'kill_long_running_queries': self.kill_long_running_queries
        }
        
        handler = action_handlers.get(action)
        if handler:
            return await handler(issue)
        else:
            print(f": {action}")
            return False
    
    async def restart_service(self, issue):
        """"""
        try:
            # 
            service_name = self.identify_affected_service(issue)
            
            # 
            await self.graceful_shutdown(service_name)
            
            # 
            await asyncio.sleep(10)
            
            # 
            await self.start_service(service_name)
            
            # 
            return await self.wait_for_service_ready(service_name, timeout=120)
            
        except Exception as e:
            print(f": {e}")
            return False
    
    async def scale_out_replicas(self, issue):
        """"""
        try:
            service_name = self.identify_affected_service(issue)
            current_replicas = await self.get_current_replica_count(service_name)
            
            # 
            new_replica_count = min(current_replicas + 2, self.config.max_replicas)
            
            await self.set_replica_count(service_name, new_replica_count)
            
            # 
            return await self.wait_for_replicas_ready(service_name, new_replica_count)
            
        except Exception as e:
            print(f": {e}")
            return False
\end{lstlisting}

\hypertarget{section-129}{%
\subsection{}\label{section-129}}

\important{}

\begin{enumerate}
\def\labelenumi{\arabic{enumi}.}
\item
  \important{}
\item
  \important{}
\item
  \important{}CI/CD
\item
  \important{}
\end{enumerate}

8''\,''

\hypertarget{section-130}{%
\subsection{}\label{section-130}}

\hypertarget{section-131}{%
\subsubsection{}\label{section-131}}

\begin{enumerate}
\def\labelenumi{\arabic{enumi}.}
\tightlist
\item
\item
\item
\end{enumerate}

\hypertarget{section-132}{%
\subsubsection{}\label{section-132}}

\begin{enumerate}
\def\labelenumi{\arabic{enumi}.}
\tightlist
\item
\item
\item
\end{enumerate}

\hypertarget{section-133}{%
\subsubsection{}\label{section-133}}

\begin{enumerate}
\def\labelenumi{\arabic{enumi}.}
\tightlist
\item
\item
\end{enumerate}

\hypertarget{section-134}{%
\subsection{}\label{section-134}}